\documentclass[UTF8,12pt,a4paper,fontset=fandol]{ctexbook}

% 支持从项目根目录或当前笔记目录编译。
\makeatletter
\def\input@path{{../}}
\makeatother
\usepackage{researchnotes}

\title{组合数学笔记}
\author{}
\date{}

\begin{document}
\maketitle
\frontmatter
\tableofcontents

\chapter{绪言与课程路线}
\label{ch:course-outline}

\section{研究对象与学习目标}

\keyterm{组合数学（combinatorics）}研究离散对象的计数、存在、构造及其结构。安排座位、分配物品、选择代表、统计路径、设计连接网络和识别对称图案，看似不同，却都需要先明确对象与限制，再寻找能够解释答案的方法。本书面向本科生，以有限集合、初等代数和基本证明方法为起点，沿着“对象建模---精确计数---结构分析---存在与构造”的路线组织内容，不把组合数学处理为孤立公式和解题技巧的汇编。

全书的学习目标是：能够把文字问题转化为无歧义的数学对象；根据约束选择分类、分步、双射、递推、生成函数或群作用等方法；独立完成主要结论的证明；区分精确公式、上下界、存在性证明与可执行算法。矩阵树定理和线性码专题需要线性代数，概率方法需要有限概率空间与期望，相关预备知识在附录中补充，不作为阅读开篇内容的障碍。

本书以定义、定理、证明和例题组成连续正文，围绕有限对象建立可核查的论证。核心结果给出证明，较长推导尽量由已学方法分解；所依赖的外部预备事实在使用处说明，不以证明思路冒充完整证明。章末习题用于检验模型、训练论证，附录给出代表性习题的解答。小规模计算能够核对例子，不能替代对任意参数成立的数学证明。

\section{全书大纲}

全书内容按知识依赖排列，各章在课程中的作用如下。
\begin{enumerate}
  \item 第\ref{ch:objects}章，绪论、组合对象与证明方法：集合语言、对象编码、基本计数原理、双射、归纳与构造。
  \item 第\ref{ch:elementary-counting}章，排列、组合与分配模型：有序与无序、允许重复与不允许重复、多重集合、圆排列和隔板法。
  \item 第\ref{ch:identities}章，二项式系数与组合恒等式：二项式与多项式定理、双重计数、带标记对象和格路径。
  \item 第\ref{ch:pigeonhole}章，抽屉原理与存在性证明：抽屉的选择、加强形式、余数模型、单调子序列和不变量。
  \item 第\ref{ch:inclusion-exclusion}章，容斥原理与二项式反演：交叠条件、错排、满射、有限上界分配与反演思想。
  \item 第\ref{ch:recurrences}章，递推关系与递归结构：组合分解、初值、线性递推、特征根和有限状态模型。
  \item 第\ref{ch:ogf}章，普通生成函数：形式幂级数、卷积、递推求解、受限分配和多变量标记。
  \item 第\ref{ch:egf}章，指数生成函数与集合分拆：标号结构、Stirling 数、Bell 数、置换循环和指数公式。
  \item 第\ref{ch:classical-structures}章，Catalan 数与整数分拆：括号、树与路径的对应，反射法、Ferrers 图及分拆恒等式。
  \item 第\ref{ch:graphs-trees}章，图、路径与树的计数：图的语言、连通性、Euler 道、树、Pr\"ufer 编码与生成树。
  \item 第\ref{ch:matching-coloring}章，匹配、染色与图的结构：Hall 定理、增广路、K\H{o}nig 定理、染色多项式和平面图入门。
  \item 第\ref{ch:symmetry}章，群作用与对称计数：轨道、稳定子、Burnside 引理、循环指标和 P\'olya 计数。
  \item 第\ref{ch:posets}章，偏序集与 M\"obius 反演：链与反链、关联代数、一般反演、Sperner 定理及格的入门。
  \item 第\ref{ch:extremal-ramsey}章，极值组合与 Ramsey 理论：极值函数、Mantel 定理、Tur\'an 图、Ramsey 数及上下界。
  \item 第\ref{ch:probabilistic}章，概率方法：随机结构、第一矩方法、随机二染色、删改法和条件期望法。
  \item 第\ref{ch:designs-codes}章，组合设计与纠错编码：区组设计、有限几何小例子、正交拉丁方、距离与编码界。
  \item 第\ref{ch:algorithms}章，组合算法与综合问题：枚举、动态规划、状态压缩、回溯，以及多种方法的比较与核验。
\end{enumerate}

\section{课程层次与知识依赖}

第\ref{ch:objects}至\ref{ch:inclusion-exclusion}章建立计数与证明基础；第\ref{ch:recurrences}至\ref{ch:classical-structures}章形成递推、生成函数和经典对象之间的联系；第\ref{ch:graphs-trees}至\ref{ch:symmetry}章转向离散结构与对称。以上十二章构成本书的主体，但矩阵树定理的完整证明、平面图的较深结论和完整 P\'olya 理论可以在首次学习时选读。第\ref{ch:posets}至\ref{ch:algorithms}章是五个拓展专题，不要求在一门短学期课程中全部讲完。

阅读时应保留以下依赖：普通生成函数依赖递推与基本计数，指数生成函数进一步依赖标号对象的分配；Catalan 数的闭式推导使用格路径与生成函数；匹配和极值图论需要图的基本语言；对称计数在章内补充有限群作用，不预设抽象代数课程；偏序反演从容斥和二项式反演出发；概率方法在使用 Ramsey 问题作为例子时衔接前一章。算法章中的基础枚举与动态规划可以提前配合相应计数章节阅读，设计与编码章则可独立选学。

基础学习路线以“每种方法能独立解决一个新问题”为验收标准，不以记住公式数量为目标。各章从具体问题引出概念和核心定理，再用直接应用、变式和反例辨明条件；章末练习覆盖概念辨析、计算建模、证明训练与综合拓展。涉及高级结果时，明确标注完整证明、证明思路或引用不证，不把未说明的结论当作前提。

\section{符号约定与参考范围}

记 $\mathbb Q,\mathbb R,\mathbb C$ 分别为有理数、实数、复数集，$\mathbb Z_{\ge0}$、$\mathbb Z_{>0}$ 分别为非负整数集和正整数集；对 $n\in\mathbb Z_{\ge0}$，记 $[n]=\{1,\ldots,n\}$，其中 $[0]=\varnothing$。有限集合 $A$ 的基数记为 $|A|$，其幂集记为 $2^A$。约定 $0!=1$，空和为 $0$，空积为 $1$；当 $n\ge0$ 而整数 $k$ 不满足 $0\le k\le n$ 时，置 $\binom nk=0$。计数问题中的参数范围由具体命题说明，不能依赖字母形状猜测。

本书区分\keyterm{标号（labeled）}与\keyterm{无标号（unlabeled）}对象，并为每个商计数问题明确等价关系。后文的 $[x^n]A(x)$ 表示形式幂级数 $A(x)$ 中 $x^n$ 的系数；图通常是有限简单无向图，需要允许重边、环或方向时另行声明。概率专题使用大写字母表示随机变量，小写字母表示其取值或固定参数。

课程范围参考 MIT 的本科组合分析课程资料\cite{mit-combinatorics}，问题组织可参考 Bogart 的开放教材\cite{bogart}，偏序、标号结构与群作用的拓展次序可与 Dartmouth 的代数组合课程资料\cite{dartmouth-combinatorics}对照。这些来源用于核对课程覆盖与提供后续阅读入口；本书的章节划分和教学安排另行组织，不复制来源中的成套习题。

\mainmatter

\chapter{绪论、组合对象与证明方法}
\label{ch:objects}

从四个人中选两名代表有 $6$ 种方案，选一名正代表和一名副代表却有 $12$ 种方案。差别不在计算技巧，而在计数对象：前者是二元素子集，后者是由两个不同元素组成的有序对。组合问题的第一步因此不是选择公式，而是确定什么算作一个对象、哪些对象合法、何时两个对象相同。

\section{计数、存在、构造与优化}

设 $X$ 是满足给定限制的有限对象集合。\keyterm{计数问题（enumeration problem）}要求 $|X|$，存在性问题要求判断 $X$ 是否非空，构造问题要求实际给出 $X$ 的一个元素；若另给代价函数 $w:X\to\mathbb R$，优化问题要求在 $X$ 中使 $w$ 最小或最大。后者只有在可行集非空时才有可行解；有限性保证非空可行集上的实值函数取得极值，却不保证容易找到极值点。

例如，每个学生有若干可选课题，且每个课题至多分给一人。可以问合法分配的数量、是否存在覆盖所有学生的分配、怎样找到分配，或怎样使总偏好得分最大。同一情境产生四个不同任务。给出一个方案只能证明存在，检查许多失败方案一般不能证明不存在；一个计数公式也不自动附带有效的构造算法。

\section{集合、映射与组合编码}

\begin{definition}
映射 $f:X\to Y$ 若满足 $f(x)=f(x')\Rightarrow x=x'$，称为\keyterm{单射（injection）}；若每个 $y\in Y$ 都有原像，称为\keyterm{满射（surjection）}；同时为单射和满射时称为\keyterm{双射（bijection）}。集合上的关系若具有自反性、对称性、传递性，称为\keyterm{等价关系（equivalence relation）}。
\end{definition}

对有限集合，单射给出 $|X|\le|Y|$，满射给出 $|X|\ge|Y|$，双射给出 $|X|=|Y|$。等价关系把集合划分成互不相交的等价类；将某些变换后的对象视为同一个方案，就是在计数等价类。构造编码时必须记录解码所需的信息，否则只能得到一个多对一映射，不能直接断言等数。

\begin{example}[子集与零一串]
对子集 $A\subseteq[n]$，令第 $i$ 位为 $1$ 当且仅当 $i\in A$，得到一个长度为 $n$ 的零一串。反向取出所有值为 $1$ 的位置，便恢复 $A$；两步互逆。因此 $[n]$ 的子集与长度为 $n$ 的零一串一一对应。此编码还保持“子集大小”等于“串中 $1$ 的个数”，可以直接处理附加限制。
\end{example}

\section{加法、乘法与等重计数原理}

\begin{theorem}[基本计数原理]
若有限集合 $X$ 是两两不交集合 $X_1,\ldots,X_m$ 的并，则 $|X|=\sum_i|X_i|$。若一个构造分成 $r$ 步，对任意合法的前 $i-1$ 步选择，第 $i$ 步都有恰好 $a_i$ 种延续，且最终选择序列与对象一一对应，则对象数为 $\prod_i a_i$。若 $f:X\to Y$ 满足每个 $y\in Y$ 恰有 $d>0$ 个原像，则 $|X|=d|Y|$。
\end{theorem}
\begin{proof}
第一式中每个元素恰在一项中被计数。第二式对步数归纳：每个长度为 $i-1$ 的合法前缀有 $a_i$ 个不同扩展，所以前缀总数逐步乘以 $a_i$。第三式将 $X$ 分成各纤维 $f^{-1}(\{y\})$，对这些互不相交的 $d$ 元集合应用第一式。
\end{proof}

前两条分别称为\keyterm{加法原理（sum rule）}和\keyterm{乘法原理（product rule）}。乘法原理不要求各步选择的具体内容互不影响，但要求相应延续数恒定；若第一步选 $x\in A$ 后有 $b(x)$ 种延续，总数应为 $\sum_{x\in A}b(x)$。例如先选一个数 $i\in[n]$，再选 $j\in[i]$，得到 $n(n+1)/2$ 个有序对，而非 $n^2$ 个。

等重条件是除法的依据。从四人中顺次选两人有 $4\cdot3$ 种，每个二元素子集恰有两种排列，故子集数为 $12/2=6$。若某些对象只有一种编码、另一些有两种编码，则不存在一个可以统一除去的“重复次数”。此外，空串与空排列各是一个对象，所以零步构造的方案数是 $1$；这与空集的元素数为 $0$ 并不矛盾。

\section{归纳、双射与最小反例}

\keyterm{数学归纳法（mathematical induction）}要求验证起点，并由 $P(n)$ 推出 $P(n+1)$。\keyterm{强归纳法（strong induction）}允许在证明 $P(n)$ 时使用所有更小且不低于起点的情形。二者可以由非负整数的良序性质解释：若存在反例，最小反例的所有较小情形成立，归纳步骤便排除了它。因此使用递归分解时，必须说明参数确实减小，并覆盖最小对象。

\begin{proposition}
对 $n\ge0$，$|2^{[n]}|=2^n$；对 $n\ge1$，偶数大小子集和奇数大小子集各有 $2^{n-1}$ 个。
\end{proposition}
\begin{proof}
子集的每个元素有选或不选两种状态，故由乘法原理得到第一式；也可按是否含 $n$ 分成两个与 $2^{[n-1]}$ 等大的类，以初值 $|2^{\varnothing}|=1$ 归纳。对第二式，固定元素 $1$，把包含 $1$ 的子集删去 $1$，把不包含 $1$ 的子集加入 $1$。此操作改变大小的奇偶性，连续做两次恢复原集合，因此给出两类之间的双射。二者和为 $2^n$，各占一半。
\end{proof}

该证明展示了两种组织方式：递推按较小对象分解，双射直接匹配两类对象。选择方法的标准是能否清楚解释对象，不是公式看起来是否简短。

\section{习题}
\begin{enumerate}
  \item\label{ex:objects-1} 从五人中选三名代表，再从代表中指定一名召集人，求方案数，并用另一种选择次序核对。
  \item\label{ex:objects-2} 设有限映射 $f:X\to Y$ 的每个纤维至多含 $d$ 个元素，证明 $|X|\le d|Y|$；给出等号成立的条件。
  \item\label{ex:objects-3} 说明从 $[n]$ 到空集的映射数，以及从空集到任意集合的映射数。区分 $n=0$ 与 $n>0$。
  \item\label{ex:objects-4} 用归纳法证明：有限集合 $A$ 的所有子集的大小之和为 $|A|2^{|A|-1}$（$|A|\ge1$），再尝试给出双重计数证明。
\end{enumerate}

\chapter{排列、组合与分配模型}
\label{ch:elementary-counting}

排列与组合公式处理的是几种标准对象。使用前先回答四个问题：元素是否可区分，顺序是否保留，是否允许重复，是否要求每个容器非空。把限制直接写进对象定义，比从题目中的“选”“排”“分”猜测公式更可靠。

\section{排列、组合与可重复选择}

\begin{definition}
从 $n$ 个不同元素中选出 $k$ 个不同元素并排成序列，称为一个 $k$-\keyterm{排列（permutation）}；只保留所选元素的集合，称为一个 $k$-\keyterm{组合（combination）}。当 $k=n$ 时，排列也称全排列。
\end{definition}

对 $0\le k\le n$，逐位选择得到排列数 $(n)_k=n(n-1)\cdots(n-k+1)=n!/(n-k)!$，其中 $(n)_0=1$。每个 $k$ 元子集恰有 $k!$ 个排列，因此
\begin{equation}
  \binom nk=\frac{(n)_k}{k!}=\frac{n!}{k!(n-k)!}.
  \label{eq:binomial-count}
\end{equation}
若允许每位重复选取，则长度为 $k$ 的序列是从位置集 $[k]$ 到元素集 $[n]$ 的映射，共有 $n^k$ 个。空位置集只有一个映射，故在这种计数语境中约定 $0^0=1$；非空集合到空集没有映射。

\begin{example}
用十个数字组成四位十进制正整数。允许重复时，首位有 $9$ 种，其余各有 $10$ 种，得 $9000$ 个；数字互不相同时，依次有 $9,9,8,7$ 种选择，得 $4536$ 个。第二位仍有 $9$ 种，是因为可以用零，但不能重复首位。这说明分步选择数可以相同，而对应的候选集合并不相同。
\end{example}

\section{多重集合排列与圆排列}

一个\keyterm{多重集合（multiset）}允许元素具有重数。若 $r$ 类元素的重数为 $n_1,\ldots,n_r\ge0$，总数为 $n$，先把所有副本临时标为不同，得到 $n!$ 个排列。去掉副本标签后，每个实际排列恰有 $n_1!\cdots n_r!$ 个原像，故排列数为
\begin{equation}
  \binom{n}{n_1,\ldots,n_r}=\frac{n!}{n_1!\cdots n_r!}.
  \label{eq:multinomial-count}
\end{equation}
例如字母串 $\mathrm{BANANA}$ 有三个 $A$、两个 $N$、一个 $B$，不同排列数为 $6!/(3!2!)=60$。

把 $n\ge1$ 个不同的人围成圆桌，只把整体旋转视为相同。固定一人的位置，其余 $n-1$ 人按顺时针方向排列，得到 $(n-1)!$ 种。固定位置是在每个等价类中选唯一代表；若再把镜像视为相同，等价关系已改变。含重复颜色时也不能一般地用线性排列数除以 $n$：例如二色交替排列可能在未转满一圈前恢复原状，各轨道未必等大。

\begin{example}
五个不同的人围成圆桌，其中甲乙必须相邻。把甲乙合成一个块后，有四个不同的圆排列单元，共 $(4-1)!$ 种，块内有 $2$ 种顺序，故共有 $12$ 种。不要求相邻时有 $24$ 种，因此甲乙不相邻也有 $12$ 种。此处块与圆桌安排可相互恢复，计数没有遗漏或重复。
\end{example}

\section{隔板法与整数向量}

将 $n\ge0$ 个相同物品放入 $k\ge1$ 个不同盒子，由各盒数量 $(x_1,\ldots,x_k)$ 唯一确定方案。它等价于非负整数方程 $x_1+\cdots+x_k=n$。

\begin{theorem}[隔板法（stars and bars）]
上述方程的非负整数解数为 $\binom{n+k-1}{k-1}$。若各变量均为正整数，则在 $n\ge k$ 时解数为 $\binom{n-1}{k-1}$，在 $n<k$ 时为零。
\end{theorem}
\begin{proof}
把 $n$ 个星号排成一行，插入 $k-1$ 个隔板，按隔板划分成 $k$ 段，段内星号数即为 $x_i$；允许隔板相邻和出现在两端。由此与含 $n$ 个星号、$k-1$ 个隔板的字符串建立双射，选隔板位置得到第一式。若 $x_i\ge1$，令 $y_i=x_i-1\ge0$，总和变成 $n-k$，应用第一式即得第二式。
\end{proof}

更一般地，若 $x_i\ge a_i$，其中 $a_i$ 为给定非负整数，先令 $y_i=x_i-a_i$；剩余总数 $n-\sum_i a_i$ 为负时无解。若同时有上界，则隔板位置受到额外约束，不能继续直接套用无约束公式。

\begin{example}
方程 $x_1+x_2+x_3=8$ 有 $\binom{10}{2}=45$ 个非负整数解；各变量均正时有 $\binom72=21$ 个；若 $x_1\ge2,x_2\ge1,x_3\ge0$，平移后总和为 $5$，有 $\binom72=21$ 个。后两者数量相同，但对应的原始解集不同。
\end{example}

\section{分配问题的统一分类}

不同物品放入不同盒子是映射问题：不限制空盒时为任意映射，每盒至多一件时为单射，每盒至少一件时为满射。相同物品放入不同盒子则是整数向量问题。若不同物品的盒子没有名称且只保留非空盒，得到\keyterm{集合分拆（set partition）}；若物品与盒子都相同，只保留各盒的正数量并按非增次序记录，得到\keyterm{整数分拆（integer partition）}。后两者将在第\ref{ch:egf}、\ref{ch:classical-structures}章讨论。

把正整数 $n$ 写成有序正整数之和称为\keyterm{整数合成（composition）}。固定 $k$ 项时有 $\binom{n-1}{k-1}$ 种；不固定项数而 $n\ge1$ 时，在一排 $n$ 个单位的 $n-1$ 个间隙中任意选择断点，故共有 $2^{n-1}$ 种。$1+3$ 与 $3+1$ 是不同合成，却是同一个整数分拆。是否保留顺序决定了模型。

\section{习题}
\begin{enumerate}
  \item\label{ex:elementary-counting-1} 求字母串 $\mathrm{MISSISSIPPI}$ 的不同排列数，并解释分母各因子的意义。
  \item\label{ex:elementary-counting-2} 六人排成一行，指定两人不相邻，有多少种排列？改为圆桌且只识别旋转后再求解。
  \item\label{ex:elementary-counting-3} 求 $x_1+x_2+x_3+x_4=12$ 且各 $x_i\ge2$ 的整数解数。
  \item\label{ex:elementary-counting-4} 把五个不同物品放入三个不同盒子，分别写出任意分配、每盒至多一件、每盒至少一件的模型；最后一种暂不要求套公式。
\end{enumerate}

\chapter{二项式系数与组合恒等式}
\label{ch:identities}

一个组合恒等式通常意味着同一对象集具有两种描述。代数计算可以验证等式，组合证明则进一步说明两侧为何应当相等。带权和也能纳入这一观点：权重往往对应为对象增加一个标记或一个附加选择。

\section{二项式定理与多项式定理}

\begin{theorem}
对非负整数 $n$ 及可交换的变量 $x,y,x_1,\ldots,x_r$，有
\[
  (x+y)^n=\sum_{k=0}^n\binom nk x^ky^{n-k},\qquad
  (x_1+\cdots+x_r)^n=
  \sum_{\substack{n_1+\cdots+n_r=n\\n_i\ge0}}
  \binom{n}{n_1,\ldots,n_r}\prod_{i=1}^r x_i^{n_i}.
\]
\end{theorem}
\begin{proof}
展开乘积时，从每个因子中选择一个加项。得到 $x^ky^{n-k}$ 恰需在 $n$ 个因子中指定 $k$ 个选择 $x$，所以系数为 $\binom nk$。一般情形先指定哪些因子选择各个 $x_i$，计数正是式\eqref{eq:multinomial-count}。
\end{proof}

这分别是\keyterm{二项式定理（binomial theorem）}与\keyterm{多项式定理（multinomial theorem）}。它们首先是有限多项式恒等式，代入数值不会产生收敛问题；非整数指数的无限展开属于另一层次，不能由有限乘积的论证直接得到。

\section{Pascal 恒等式与边界条件}

设 $n\ge1$。按一个 $k$ 元子集是否包含指定元素 $n$ 分类，得到
\[
  \binom nk=\binom{n-1}{k-1}+\binom{n-1}{k}.
\]
不包含 $n$ 时从 $[n-1]$ 选 $k$ 个，包含时还须选 $k-1$ 个。配合越界下标取零及 $\binom00=1$，可逐行生成 Pascal 三角形。取补集给出 $\binom nk=\binom n{n-k}$；两式分别来自分类与双射，不必先依赖阶乘运算。

对全部子集按大小分类，得 $\sum_{k=0}^n\binom nk=2^n$。当 $n\ge1$ 时，偶数大小与奇数大小子集等数，故 $\sum_{k=0}^n(-1)^k\binom nk=0$；$n=0$ 时该和为 $1$。区分这个边界是后续反演证明的关键。

\section{双重计数与带标记对象}

\keyterm{双重计数（double counting）}是对同一个有限集合按不同属性求和。令
\[
  \mathcal I=\{(A,a):A\subseteq[n],\ a\in A\}.
\]
先选大小为 $k$ 的 $A$，再从中标记一个元素，得到 $\sum_k k\binom nk$；先选标记 $a$，再任意选择其余 $n-1$ 个元素是否入选，得到 $n2^{n-1}$。因此对 $n\ge1$，
\begin{equation}
  \sum_{k=0}^n k\binom nk=n2^{n-1}.
  \label{eq:marked-subsets}
\end{equation}
同样地，标记一个有序的不同元素对，得到 $\sum_k k(k-1)\binom nk=n(n-1)2^{n-2}$（$n\ge2$）。由于 $k^2=k(k-1)+k$，可据此求出平方权重和。

\begin{example}
从八人中任意选一个非空委员会并指定一名负责人。先定负责人有 $8$ 种，其余七人各可入选或不入选，共 $8\cdot2^7=1024$ 种。若先定委员会大小，则得到 $\sum_{k=1}^8 k\binom8k$。这里负责人由委员会内部产生，不能把所有委员会数直接乘以 $8$，因为并非每个人都在每个委员会中。
\end{example}

\section{Vandermonde 恒等式与格路径}

\begin{theorem}[Vandermonde 恒等式]
对非负整数 $r,s,m$，有
\begin{equation}
  \sum_{k=0}^{m}\binom rk\binom{s}{m-k}=\binom{r+s}{m}.
  \label{eq:vandermonde}
\end{equation}
\end{theorem}
\begin{proof}
把 $r+s$ 个不同元素分成大小为 $r,s$ 的两个不交集合。一个 $m$ 元子集从第一部分选出唯一确定的 $k$ 个元素，从第二部分选出 $m-k$ 个元素。对 $k$ 分类即得；不可能的 $k$ 对应零项，因而无需另改求和范围。
\end{proof}

取 $r=s=n,m=n$ 并使用对称性，得 $\sum_{k=0}^n\binom nk^2=\binom{2n}{n}$。另一个常用恒等式为 $\sum_{j=k}^n\binom jk=\binom{n+1}{k+1}$（$0\le k\le n$）：计数 $[n+1]$ 的 $(k+1)$ 元子集，若最大元素为 $j+1$，其余 $k$ 个从 $[j]$ 中选取。

\keyterm{格路径（lattice path）}在本章指每步只向右或向上移动一个单位的路径。从 $(0,0)$ 到 $(a,b)$，其中 $a,b\ge0$，共有 $a+b$ 步，其中 $a$ 步向右，故路径数为 $\binom{a+b}{a}$。若必须经过 $(u,v)$ 且 $0\le u\le a,0\le v\le b$，分成经过该点前后的两段，得到
\[
  \binom{u+v}{u}\binom{a+b-u-v}{a-u}.
\]
例如从 $(0,0)$ 到 $(4,3)$ 有 $35$ 条路径，其中经过 $(2,1)$ 的有 $3\cdot6=18$ 条，因此避开该点的有 $17$ 条。多个禁点的交叠条件将引向容斥原理。

\section{习题}
\begin{enumerate}
  \item\label{ex:identities-1} 用双重计数证明 $\binom nk\binom kr=\binom nr\binom{n-r}{k-r}$，其中 $0\le r\le k\le n$。
  \item\label{ex:identities-2} 求 $\sum_{k=0}^n k^2\binom nk$，分别处理 $n=0,1$ 后写出适用的一般公式。
  \item\label{ex:identities-3} 从 $(0,0)$ 到 $(5,4)$ 的格路径中，同时经过 $(1,1)$ 和 $(3,2)$ 的有多少条？若两个必经点不能按坐标偏序排列会怎样？
  \item\label{ex:identities-4} 给出“曲棍球杆”恒等式的递推证明，并与按最大元素分类的证明比较。
\end{enumerate}

\chapter{抽屉原理与存在性证明}
\label{ch:pigeonhole}

计数也能用于证明某种结构必然存在。如果对象太多，而在不出现目标结构时每类只能容纳很少对象，就会产生矛盾。这里的困难通常不是求和，而是设计恰当的分类。

\section{基本形式与容量形式}

\begin{theorem}[抽屉原理（pigeonhole principle）]
将 $N\ge0$ 个对象放入 $k\ge1$ 个盒子，至少一个盒子含有不少于 $\lceil N/k\rceil$ 个对象。更一般地，若第 $i$ 盒的容量为非负整数 $c_i$，且 $N>\sum_{i=1}^k c_i$，则至少一盒超过其容量。
\end{theorem}
\begin{proof}
若每盒均不超过容量，总数至多为容量之和，得一般形式。对基本形式，若每盒都少于 $\lceil N/k\rceil$，每盒对象数都严格小于 $N/k$，总数便小于 $N$，矛盾。
\end{proof}

写 $N=qk+r$，其中 $0\le r<k$。使 $r$ 个盒子各有 $q+1$ 个对象，其余各有 $q$ 个，就达到上述下界。因此在没有附加结构时，不能普遍保证更大的占用数。证明阈值最优通常需要两部分：超过阈值必然发生，以及在阈值以下确实存在避免目标的构造。

\section{余数、区间与部分和}

两个整数除以 $m\ge1$ 的余数相同，当且仅当其差被 $m$ 整除。因此从 $m+1$ 个整数中必能找到两个差被 $m$ 整除。若问题涉及连续区间的和，应先把区间和改写为部分和之差。

\begin{proposition}
对任意 $n\ge1$ 个整数 $a_1,\ldots,a_n$，存在 $1\le i\le j\le n$，使 $a_i+\cdots+a_j$ 被 $n$ 整除。
\end{proposition}
\begin{proof}
置 $s_0=0$、$s_j=\sum_{t=1}^j a_t$。$n+1$ 个部分和只有 $n$ 种可能余数，故有 $0\le u<v\le n$ 满足 $s_u\equiv s_v\pmod n$。于是非空区间之和 $a_{u+1}+\cdots+a_v=s_v-s_u$ 被 $n$ 整除。
\end{proof}

类似地，把 $[0,1]$ 划成 $m$ 个长度为 $1/m$ 的区间，采用左闭右开、最后一段包含右端点的约定，则任意 $m+1$ 个点中至少两个落在同一段，距离不超过 $1/m$。端点约定保证每个点只属于一个盒子；它是证明的一部分。

\section{Erd\H{o}s--Szekeres 单调子序列定理}

\begin{theorem}
设 $r,s\ge2$。任意由 $(r-1)(s-1)+1$ 个互不相同实数组成的序列，含有长度为 $r$ 的严格递增子序列或长度为 $s$ 的严格递减子序列。
\end{theorem}
\begin{proof}
对位置 $i$，令 $p_i,q_i$ 分别为以该位置结尾的最长严格递增、严格递减子序列的长度。若结论不成立，则所有 $(p_i,q_i)$ 落在 $[r-1]\times[s-1]$ 中。对 $i<j$，若 $a_i<a_j$，把 $a_j$ 接到以 $i$ 结尾的递增子序列后，得 $p_j\ge p_i+1$；若 $a_i>a_j$，同理 $q_j\ge q_i+1$。故各有序对互不相同，但位置数超过可用有序对数，矛盾。
\end{proof}

子序列只要求下标递增，不要求相邻。阈值不能降低：把 $1,\ldots,(r-1)(s-1)$ 分成 $r-1$ 个连续数值块，每块含 $s-1$ 个数；块按数值从小到大排列，块内递减。递增子序列每块至多取一个，递减子序列只能在一块内，因而两种目标均不存在。

\section{极端原理、不变量与构造}

\keyterm{极端原理（extremal principle）}通过选择某个量最大的、最小的或最短的对象，利用其不可进一步改进的性质。它要求极端对象确实存在：有限非空集合自动满足这一点，一般无限集合则未必如此。\keyterm{不变量（invariant）}是在每次合法操作后保持不变的量；若初态与目标态的不变量不同，就证明目标不可达。

\begin{example}[棋盘染色]
从 $8\times8$ 棋盘删除两个对角顶角后，不能用 $1\times2$ 骨牌覆盖。将棋盘交替染成黑白两色，每块骨牌必覆盖一黑一白；两个对角顶角同色，删除后两种颜色的格数分别为 $30$ 与 $32$，而任何骨牌覆盖要求二者相等，故不可能。这里证明的是不存在，并未依靠逐一枚举铺法。
\end{example}

\begin{example}[同一集合的两种抽屉]
从 $[2n]$ 中选出 $n+1$ 个数。按 $\{1,2\},\{3,4\},\ldots,\{2n-1,2n\}$ 分组，某组取了两个数，故有相邻数。若要证明一个整除另一个，改为按奇数部分分组：每个正整数唯一写成 $2^a u$，其中 $u$ 为奇数；$[2n]$ 中只有 $n$ 种可能的奇数部分。同组的两个数为 $2^a u,2^b u$，较小者整除较大者。不同目标要求不同抽屉。
\end{example}

\section{习题}
\begin{enumerate}
  \item\label{ex:pigeonhole-1} 证明任意六个整数中有两个之差被 $5$ 整除，并构造五个整数使任意两数之差都不被 $5$ 整除。
  \item\label{ex:pigeonhole-2} 在单位正方形内取五个点，证明有两点距离不超过 $\sqrt2/2$；说明边界点怎样归类。
  \item\label{ex:pigeonhole-3} 构造长度为 $9$ 的不同整数序列，使其没有长度为 $4$ 的递增或递减子序列。
  \item\label{ex:pigeonhole-4} 初始写有整数 $1,2,\ldots,10$，每次擦掉两个数 $a,b$ 并写入 $a+b$。证明最后一个数不依赖操作次序；若改为 $a+b+1$，最后结果是多少？
\end{enumerate}

\chapter{容斥原理与二项式反演}
\label{ch:inclusion-exclusion}

若多个坏条件可以同时出现，分别扣除它们会重复扣除交集。容斥的作用是精确修正每个对象的总贡献，而不是凭记忆交替添加若干项。使用时先确定全集，再定义坏事件及其交集。

\section{容斥公式及其逐元素证明}

\begin{theorem}[容斥原理（principle of inclusion--exclusion）]
设 $\Omega$ 为有限集合，$A_1,\ldots,A_m\subseteq\Omega$，则
\begin{equation}
  \left|\Omega\setminus\bigcup_{i=1}^m A_i\right|
  =\sum_{S\subseteq[m]}(-1)^{|S|}
    \left|\bigcap_{i\in S}A_i\right|,
  \label{eq:inclusion-exclusion}
\end{equation}
其中空指标交集定义为 $\Omega$。
\end{theorem}
\begin{proof}
固定 $\omega\in\Omega$，假设它恰好属于 $t$ 个坏事件。右侧包含它的项对应这 $t$ 个事件的所有子集，因此其总贡献为 $\sum_{j=0}^t(-1)^j\binom tj$。当 $t=0$ 时为 $1$，当 $t>0$ 时为零，恰好与左侧的计数一致。对所有 $\omega$ 求和即得。
\end{proof}

记事件的指示函数为 $\mathbf1_A$，公式也可由 $\mathbf1_{\Omega\setminus\cup_iA_i}=\prod_i(1-\mathbf1_{A_i})$ 展开后求和得到。若只需求至少满足一个坏条件的数量，用 $|\Omega|$ 减去式\eqref{eq:inclusion-exclusion}即可；若各坏事件本来互斥，则所有二重及以上交集为空，公式退化为加法原理。

\section{错排与受限位置}

\begin{definition}
$[n]$ 上的排列 $\pi$ 若对每个 $i$ 都满足 $\pi(i)\ne i$，称为\keyterm{错排（derangement）}。错排数记为 $D_n$，并规定 $D_0=1$。
\end{definition}

在全部 $n!$ 个排列中取坏事件 $A_i=\{\pi:\pi(i)=i\}$。指定 $j$ 个位置固定后，其余位置任意排列，交集大小为 $(n-j)!$；故
\begin{equation}
  D_n=\sum_{j=0}^n(-1)^j\binom nj(n-j)!
      =n!\sum_{j=0}^n\frac{(-1)^j}{j!}.
  \label{eq:derangements}
\end{equation}
例如 $D_4=24-24+12-4+1=9$。若要求恰有 $r$ 个固定点，先选这些位置，再使其余位置错排，得到 $\binom nrD_{n-r}$；若只要求至少 $r$ 个固定点，不能使用同一公式。

当 $n\ge2$ 时还有 $D_n=(n-1)(D_{n-1}+D_{n-2})$。固定 $\pi(1)=j\ne1$。若 $\pi(j)=1$，删除这个二循环，剩 $D_{n-2}$ 种；若 $\pi(j)\ne1$，设 $\pi(k)=1$，删除 $1$ 并改令 $k$ 指向 $j$，得到其余 $n-1$ 个元素上的错排。逆向在指向 $j$ 的箭头中插入 $1$，说明后一对应可逆。再乘 $j$ 的 $n-1$ 种选择即得。

一般禁位条件可放在 $n\times n$ 棋盘的格集 $B$ 上，格 $(i,j)$ 表示禁止 $\pi(i)=j$。令 $r_k$ 为在 $B$ 中放置 $k$ 个互不同行且互不同列的车的方法数，$r_0=1$，则合法排列数为 $\sum_{k=0}^n(-1)^k r_k(n-k)!$。这是对禁位事件容斥：只有互不同行列的事件交集非空，指定 $k$ 个像后还剩 $(n-k)!$ 种。多项式 $\sum_k r_kx^k$ 称为\keyterm{车多项式（rook polynomial）}。

\section{满射与有上界的分配}

从 $[n]$ 到 $[k]$ 的满射要求每个像值出现。设 $n\ge0,k\ge1$，令 $A_i$ 表示像中没有 $i$；禁止 $j$ 个像值后有 $(k-j)^n$ 个映射，因此满射数为
\[
  \sum_{j=0}^k(-1)^j\binom kj(k-j)^n.
\]
沿用空映射约定 $0^0=1$。例如五个不同物品放进三个不同非空盒子，方案数为 $3^5-3\cdot2^5+3=150$。

对非负整数上界 $u_1,\ldots,u_k$，考虑 $x_1+\cdots+x_k=n$ 且 $0\le x_i\le u_i$。定义
\[
  H_k(t)=\begin{cases}\binom{t+k-1}{k-1},&t\ge0,\\0,&t<0.\end{cases}
\]
若指标集 $S$ 中的变量违反上界，令 $y_i=x_i-(u_i+1)$（$i\in S$），其余不变，交集计数为 $H_k(n-\sum_{i\in S}(u_i+1))$。故解数等于
\begin{equation}
  \sum_{S\subseteq[k]}(-1)^{|S|}
  H_k\left(n-\sum_{i\in S}(u_i+1)\right).
  \label{eq:bounded-compositions}
\end{equation}
例如 $x_1+x_2+x_3=8$、$0\le x_i\le4$ 时，总数为 $\binom{10}{2}-3\binom52=15$；两变量同时至少为 $5$ 已不可能，所以没有后续非零项。

\section{二项式反演及其适用形式}

\begin{theorem}[二项式反演（binomial inversion）]
对实数或复数序列 $(a_n)_{n\ge0},(b_n)_{n\ge0}$，有
\begin{equation}
  b_n=\sum_{k=0}^n\binom nk a_k
  \quad\Longleftrightarrow\quad
  a_n=\sum_{k=0}^n(-1)^{n-k}\binom nk b_k.
  \label{eq:binomial-inversion}
\end{equation}
\end{theorem}
\begin{proof}
代入第一式并交换有限求和，反演式右侧变为
\[
 \sum_{j=0}^n a_j\binom nj
 \sum_{k=j}^n(-1)^{n-k}\binom{n-j}{k-j}.
\]
内层和在 $j<n$ 时为零，在 $j=n$ 时为 $1$，故结果为 $a_n$。反向代入得到同样的二项式交错和，或利用第一式作为对角元为 $1$ 的三角递推逐项唯一确定 $a_n$，从而得到等价性。
\end{proof}

例如把排列按非固定点的集合分类，得 $n!=\sum_{k=0}^n\binom nkD_k$。二项式反演立即恢复错排公式。这里的变换把“恰好使用某个支撑集”的计数汇总为“支撑集任意”的计数，反演则将二者分离。

\section{习题}
\begin{enumerate}
  \item\label{ex:inclusion-exclusion-1} 求六个不同对象的错排数，以及恰有两个固定点的排列数。
  \item\label{ex:inclusion-exclusion-2} 求长度为 $5$、使用字母 $A,B,C$ 且每种字母都出现的字符串数。
  \item\label{ex:inclusion-exclusion-3} 求 $x_1+x_2+x_3+x_4=10$ 且 $0\le x_i\le3$ 的整数解数。
  \item\label{ex:inclusion-exclusion-4} 若 $b_n=2^n$ 且 $b_n=\sum_{k=0}^n\binom nk a_k$，用反演求 $a_n$，再解释其组合意义。
\end{enumerate}

\chapter{递推关系与递归结构}
\label{ch:recurrences}

当大对象可以唯一地分解成较小对象时，计数自然产生\keyterm{递推关系（recurrence relation）}。建立递推时需先证明分解，求解递推时需另外处理初值。一个猜出的数列即使符合若干小值，也不能替代这两项工作。

\section{从对象分解建立递推}

用长度为 $1$ 和 $2$ 的砖覆盖 $1\times n$ 棋盘，记方案数为 $a_n$。空棋盘有一种空铺法，故 $a_0=1$；$a_1=1$。对 $n\ge2$，最后一块砖长度为 $1$ 或 $2$，删除后分别得到长度为 $n-1$ 或 $n-2$ 的铺法，而且可唯一补回。因此
\begin{equation}
  a_n=a_{n-1}+a_{n-2},\qquad a_0=a_1=1.
  \label{eq:tiling-recurrence}
\end{equation}
前几项为 $1,1,2,3,5,8$。若 $F_0=0,F_1=1,F_n=F_{n-1}+F_{n-2}$，则 $a_n=F_{n+1}$。

长度为 $n$ 且不含连续两个 $1$ 的零一串，记其数目为 $b_n$。当 $n\ge2$ 时，末位为 $0$ 的串由任意合法的 $n-1$ 位串补 $0$ 得到；末位为 $1$ 时末两位必须是 $01$，由合法的 $n-2$ 位串补 $01$ 得到。故同样有 $b_n=b_{n-1}+b_{n-2}$，但 $b_0=1,b_1=2$，所以 $b_n=F_{n+2}$。相同递推式不代表相同计数序列。

\section{常系数线性递推与特征方程}

设 $r\ge1$，$c_1,\ldots,c_r$ 为常数且 $c_r\ne0$，考虑
\[
  a_n=c_1a_{n-1}+\cdots+c_ra_{n-r}\quad(n\ge r).
\]
指定 $a_0,\ldots,a_{r-1}$ 后，每一后项均被唯一确定。代入 $a_n=\lambda^n$，得到\keyterm{特征方程（characteristic equation）} $\lambda^r-c_1\lambda^{r-1}-\cdots-c_r=0$。若有 $r$ 个不同的复根 $\lambda_1,\ldots,\lambda_r$，各指数序列都是解；前 $r$ 项组成的 Vandermonde 矩阵可逆，所以可唯一选择系数 $\alpha_j$ 匹配任意初值。唯一性给出 $a_n=\sum_j\alpha_j\lambda_j^n$。

若某根 $\lambda$ 的重数为 $m$，应使用 $\lambda^n,n\lambda^n,\ldots,n^{m-1}\lambda^n$；一般形式由第\ref{ch:ogf}章的部分分式展开推出，不能仍只保留一个指数项。实系数与实初值保证共轭根的系数也配成共轭，因此最终结果为实数。

\begin{example}
式\eqref{eq:tiling-recurrence}的特征根为 $\phi=(1+\sqrt5)/2$ 与 $\psi=(1-\sqrt5)/2$，匹配初值得
\[
  a_n=\frac{\phi^{n+1}-\psi^{n+1}}{\sqrt5}.
\]
它虽然含无理数，却由递推与整数初值保证取整数。对 $u_n=2u_{n-1}-u_{n-2}$、$u_0=0,u_1=1$，特征根 $1$ 为二重根，实际解为 $u_n=n$，展示了重根多项式因子的必要性。
\end{example}

\section{非齐次递推与累积量}

对 $a_n=ca_{n-1}+f(n)$（$n\ge1$），反复代入并归纳可得
\begin{equation}
  a_n=c^na_0+\sum_{j=1}^n c^{n-j}f(j).
  \label{eq:inhomogeneous-recurrence}
\end{equation}
这个公式直接包括常数、多项式和指数驱动项。若使用待定系数法，应将候选特解代回验证；驱动项与齐次解重复时通常需要额外乘以 $n$，式\eqref{eq:inhomogeneous-recurrence}可以解释这种累积效应。

\begin{example}[Hanoi 塔的最少移动次数]
三个柱上有 $n$ 个大小不同的盘，每次移动最上面的一个，且大盘不能压在小盘上。记将完整盘堆移到指定另一柱的最少步数为 $T_n$，$T_0=0$。任何成功方案在最大盘第一次移动前，必须把上面的 $n-1$ 个盘移到第三柱；在最大盘最后一次移入目标柱后，又须把其余 $n-1$ 个盘移到它上面。这两段各至少需要 $T_{n-1}$ 步，最大盘至少移动一次，故 $T_n\ge2T_{n-1}+1$。先搬小盘、搬最大盘、再搬回小盘的递归方案达到此界，所以 $T_n=2T_{n-1}+1=2^n-1$。递推等式在此依赖下界证明和达到下界的构造。
\end{example}

\section{多状态递推与转移矩阵}

若未来合法性依赖部分历史，就应把这些历史信息保留在状态中。对避免 $11$ 的零一串，令 $u_n$ 计数末位为 $0$ 的串，$v_n$ 计数末位为 $1$ 的串；把空串归入初始的零状态，取 $u_0=1,v_0=0$。则
\[
 \begin{pmatrix}u_{n+1}\\v_{n+1}\end{pmatrix}
 =\begin{pmatrix}1&1\\1&0\end{pmatrix}
  \begin{pmatrix}u_n\\v_n\end{pmatrix},\qquad b_n=u_n+v_n.
\]
矩阵第 $i$ 行第 $j$ 列表示从旧状态 $j$ 转移到新状态 $i$ 的选择数。一般有限状态计数可写成 $w_{n+1}=Mw_n$，解为 $w_n=M^nw_0$，最终再对接受状态求和。状态必须足以判断下一步是否合法，但无需记录对未来无影响的全部历史。

\section{习题}
\begin{enumerate}
  \item\label{ex:recurrences-1} 求用长度 $1,2$ 的砖覆盖长度为 $8$ 的棋盘的方案数，并说明为何空棋盘初值必须为 $1$。
  \item\label{ex:recurrences-2} 求解 $a_n=3a_{n-1}+2$、$a_0=1$，用代回与初值两项核对答案。
  \item\label{ex:recurrences-3} 建立不含连续三个 $1$ 的长度为 $n$ 的零一串的递推与初值，计算 $n=5$ 时的数量。
  \item\label{ex:recurrences-4} 对有重根的递推 $a_n=4a_{n-1}-4a_{n-2}$、$a_0=1,a_1=4$，求闭式并验证。
\end{enumerate}

\chapter{普通生成函数}
\label{ch:ogf}

生成函数把一个序列的全部信息存入各次幂的系数，使递推和组合分解转化为代数运算。本章使用形式幂级数，不把 $x$ 看成必须代入数值的实变量。因此运算的合法性来自逐系数的有限性，而非解析收敛。

\section{形式幂级数与系数提取}

\begin{definition}
对有理数序列 $(a_n)_{n\ge0}$，形式表达式 $A(x)=\sum_{n\ge0}a_nx^n$ 称为其\keyterm{普通生成函数（ordinary generating function）}，记 $[x^n]A(x)=a_n$。所有这类表达式组成 $\mathbb Q[[x]]$；相等定义为各系数分别相等，加法逐项进行，乘法定义为
\[
 [x^n](A(x)B(x))=\sum_{k=0}^n a_kb_{n-k}.
\]
\end{definition}

每个乘积系数只涉及有限项，故乘法良定义，并满足通常的结合、交换和分配律。若 $a_0\ne0$，方程 $A(x)B(x)=1$ 逐项给出 $b_0=1/a_0$ 与 $b_n=-a_0^{-1}\sum_{k=1}^n a_kb_{n-k}$，所以乘法逆唯一存在；若 $a_0=0$ 则不可能有逆。特别地，$1/(1-x)=\sum_{n\ge0}x^n$ 是形式恒等式。

当 $B(0)=0$ 时，复合 $A(B(x))=\sum_{k\ge0}a_kB(x)^k$ 也良定义，因为 $B(x)^k$ 的最低可能次数至少为 $k$。若内层常数项非零，则同一次幂可能收到无限项贡献，不能自动作形式复合。系数可按需要扩展到 $\mathbb R$ 或 $\mathbb C$，以处理特征根；组合计数序列本身仍取整数。

\section{加法、乘法、移位与标记}

若两类对象互斥且大小定义相同，其并集对应生成函数相加。若对象唯一分为一对有序组件，大小分别为 $k,n-k$，且两组件可独立选择，则整体数为 $\sum_{k=0}^n a_kb_{n-k}$，对应生成函数相乘。这一运算称为\keyterm{Cauchy 卷积（Cauchy product）}。若组件共享同一批不同标签，还需先分配标签，普通卷积便不够，见第\ref{ch:egf}章。

移位必须处理起点。例如 $\sum_{n\ge1}a_{n-1}x^n=xA(x)$，而 $\sum_{n\ge0}a_{n+1}x^n=(A(x)-a_0)/x$。形式微分定义为 $A'(x)=\sum_{n\ge1}na_nx^{n-1}$，故 $xA'(x)$ 对应在每个大小为 $n$ 的对象中标记一个位置；它满足乘法法则，可以逐系数验证。常用展开
\begin{equation}
  \frac{1}{(1-x)^k}=\sum_{n\ge0}\binom{n+k-1}{k-1}x^n\quad(k\ge1)
  \label{eq:negative-binomial}
\end{equation}
来自 $k$ 个几何级数相乘，系数恰是隔板法计数。

\section{从递推到有理生成函数}

对式\eqref{eq:tiling-recurrence}，令 $A(x)=\sum_{n\ge0}a_nx^n$，将 $n\ge2$ 的递推乘以 $x^n$ 后求和，得到
\[
 A(x)-1-x=x(A(x)-1)+x^2A(x),\qquad
 A(x)=\frac{1}{1-x-x^2}.
\]
这里每个等式均逐系数成立。分母分解为 $(1-\phi x)(1-\psi x)$，部分分式展开配合几何级数可恢复上一章的闭式。

\begin{proposition}
生成函数 $A(x)$ 可写成 $P(x)/Q(x)$，其中 $P,Q$ 为多项式且 $Q(0)=1$，当且仅当其系数从某一项起满足常系数线性递推。
\end{proposition}
\begin{proof}
写 $Q(x)=1-c_1x-\cdots-c_rx^r$，比较 $Q(x)A(x)=P(x)$ 中超过 $\deg P$ 且至少为 $r$ 的系数，得到递推。反之，将从某项起成立的递推移到一边，乘以 $x^n$ 求和；除有限个初始系数外均抵消，所以 $Q(x)A(x)$ 是多项式。
\end{proof}

若 $r$ 阶递推从 $n=r$ 起成立且 $c_r\ne0$，所得分子次数小于 $r$。在复数域分解 $Q(x)=\prod_j(1-\lambda_jx)^{m_j}$，其中不同 $\lambda_j$ 为非零特征根。部分分式没有多项式部分，而
\[
 [x^n](1-\lambda x)^{-h}=\binom{n+h-1}{h-1}\lambda^n.
\]
于是每个重数为 $m_j$ 的根贡献一个次数至多 $m_j-1$ 的多项式乘 $\lambda_j^n$。这证明了上一章给出的重根形式；分式恒等式的唯一性也保证这些项足以且唯一地表达由初值确定的序列。

\section{分配、零钱兑换与多变量生成函数}

变量 $x_i$ 允许取 $0,\ldots,u_i$ 时，其大小生成函数为 $1+x+\cdots+x^{u_i}$，所以总和为 $n$ 的解数是
\[
 [x^n]\prod_{i=1}^k(1+x+\cdots+x^{u_i})
 =[x^n]\frac{\prod_i(1-x^{u_i+1})}{(1-x)^k}.
\]
展开分子并用式\eqref{eq:negative-binomial}，正好得到容斥式\eqref{eq:bounded-compositions}。两种方法不同，但计数的是同一组整数向量。

\begin{example}[零钱兑换]
用面值 $1,2,5$ 的硬币兑换 $10$，每种不限数量，同面值硬币不可区分。生成函数为 $1/((1-x)(1-x^2)(1-x^5))$。固定五元硬币数为 $c=0,1,2$，二元硬币数分别有 $6,3,1$ 种，一元硬币数随即确定，故共有 $10$ 种。若把投币先后次序也视为方案的一部分，模型已改为字符串，不能沿用这个乘积。
\end{example}

再以 $y$ 记录硬币枚数，生成函数变为 $1/((1-yx)(1-yx^2)(1-yx^5))$，$[x^{10}y^4]$ 计数恰用四枚凑成 $10$ 的方案。每种面值为正数，保证固定总金额只涉及有限个配置。一般多变量生成函数应明确每个变量标记的统计量，提取系数后再将其他变量赋值。

\section{习题}
\begin{enumerate}
  \item\label{ex:ogf-1} 写出 $x_1+x_2+x_3=n$ 中 $x_1$ 为非负偶数、$x_2\ge1$、$0\le x_3\le2$ 的普通生成函数，并求 $n=5$ 的解数。
  \item\label{ex:ogf-2} 从 $a_n=2a_{n-1}+1,a_0=0$ 推导生成函数并提取系数。
  \item\label{ex:ogf-3} 用面值 $1,2,5$ 的硬币恰好四枚凑成 $10$，求方案数；写出相应双变量系数表达式。
  \item\label{ex:ogf-4} 形式等式 $1/(1-x)=\sum_{n\ge0}x^n$ 是否允许直接令 $x=2$ 并把右侧理解为通常的数项级数？解释原因。
\end{enumerate}

\chapter{指数生成函数与集合分拆}
\label{ch:egf}

不同标签会产生额外的分配因子。例如把 $n$ 个人分给两个有名称的队伍，大小分别为 $k,n-k$，必须先选择哪些人进入第一队。指数生成函数把这样的二项式因子吸收到分母中，使标号结构的组合仍能使用乘法。

\section{标签分配与指数型乘法}

\begin{definition}
若 $a_n$ 计数固定标签集 $[n]$ 上的一类结构，其\keyterm{指数生成函数（exponential generating function）}为 $A(x)=\sum_{n\ge0}a_nx^n/n!$。相同大小的标签集通过重标号携带相同类型结构，且重标号与复合相容。
\end{definition}

设一个结构由两个有序组件组成，标签集互不相交且并为 $[n]$。先选第一组件的 $k$ 个标签，再分别赋予结构，得
\begin{equation}
 c_n=\sum_{k=0}^n\binom nk a_kb_{n-k},\qquad
 \sum_{n\ge0}c_n\frac{x^n}{n!}=A(x)B(x).
 \label{eq:egf-product}
\end{equation}
等式也可由乘积展开逐项验证，因为 $1/(k!(n-k)!)=\binom nk/n!$。注意 $[x^n]A(x)=a_n/n!$，实际对象数还须乘以 $n!$。普通生成函数和指数生成函数可编码同一数列，但其乘法的组合含义不同。

\section{第二类 Stirling 数与 Bell 数}

\begin{definition}
$[n]$ 的\keyterm{集合分拆（set partition）}是若干非空、两两不交且并为 $[n]$ 的子集所组成的集合，各子集称为块。恰有 $k$ 个块的数量称为第二类 Stirling 数，记为 $S(n,k)$；规定 $S(0,0)=1$，$n>0$ 时 $S(n,0)=0$，其余不可能的参数取零。
\end{definition}

对 $n\ge1$，元素 $n$ 所在块若是单元素块，删除后剩 $S(n-1,k-1)$ 种；否则从该块删除 $n$，得到 $k$ 块分拆，再从 $k$ 个块中选择插入位置。故
\[
 S(n,k)=S(n-1,k-1)+kS(n-1,k).
\]
将 $k$ 个无名称块分别命名为 $1,\ldots,k$，得到满射的纤维，因此满射数为 $k!S(n,k)$。结合容斥，得到 $S(n,k)=\frac1{k!}\sum_{j=0}^k(-1)^j\binom kj(k-j)^n$（$k\ge1$）。

\keyterm{Bell 数（Bell number）} $B_n=\sum_{k=0}^nS(n,k)$ 计数全部集合分拆，$B_0=1$。在 $[n+1]$ 的分拆中，选择与 $n+1$ 同块的 $n-k$ 个元素，余下 $k$ 个任意分拆，得到 $B_{n+1}=\sum_{k=0}^n\binom nkB_k$。

\begin{example}
四个人分成两组且组无名称，有 $S(4,2)=7$ 种：大小为 $1+3$ 的有 $4$ 种，大小为 $2+2$ 的有 $\binom42/2=3$ 种。分给两个有名称且均非空的实验室，则有 $2!\cdot7=14$ 种。四人的全部分组数为 $B_4=1+7+6+1=15$。
\end{example}

\section{第一类 Stirling 数与置换循环}

一个置换可以唯一分解为不交循环，循环内部的书写起点与循环之间的书写顺序不属于额外结构。令 $c(n,k)$ 计数 $[n]$ 上恰有 $k$ 个循环的置换，称为无符号第一类 Stirling 数；取 $c(0,0)=1$，其他不可能参数为零。

元素 $n$ 或者单独形成一个循环，或者插入某个已有元素之后。前者贡献 $c(n-1,k-1)$，后者有 $n-1$ 个插入位置，删除 $n$ 能唯一恢复原置换。因此对 $n\ge1$，
\[
 c(n,k)=c(n-1,k-1)+(n-1)c(n-1,k),\qquad
 \sum_{k=0}^n c(n,k)t^k=t(t+1)\cdots(t+n-1).
\]
多项式恒等式由递推归纳，$n=0$ 时右侧为空积 $1$。有符号约定另记为 $s(n,k)=(-1)^{n-k}c(n,k)$，不能把它与这里的非负计数混用。例如 $c(n,1)=(n-1)!$（$n\ge1$），因为一个循环可固定最小元素为书写起点，其余任意排列。

\section{指数公式与非空组件}

\begin{theorem}[标号结构的指数公式]
设组件均非空，其指数生成函数为 $C(x)$，故 $C(0)=0$。若整体结构可唯一分解为一组无序、标签互不相交的组件，则恰有 $k$ 个组件的生成函数为 $C(x)^k/k!$，全部结构的生成函数为 $\exp(C(x))$。
\end{theorem}
\begin{proof}
由式\eqref{eq:egf-product}反复相乘，$C(x)^k$ 计数有顺序的 $k$ 个组件。组件的标签集非空且互不相交，故它们作为带标签对象彼此不同；忘掉顺序后，每个无序组件集恰有 $k!$ 个原像，可以等重相除。对 $k$ 求和得到指数函数。固定总大小 $n$ 时至多有 $n$ 个非空组件，因而逐系数求和有限。
\end{proof}

一个非空标签集本身恰有一种块结构，所以 $C(x)=e^x-1$，得到
\begin{equation}
 \sum_{n\ge0}S(n,k)\frac{x^n}{n!}=\frac{(e^x-1)^k}{k!},\qquad
 \sum_{n\ge0}B_n\frac{x^n}{n!}=\exp(e^x-1).
 \label{eq:stirling-bell-egf}
\end{equation}
一个大小为 $j\ge1$ 的置换循环有 $(j-1)!$ 种，因此组件生成函数为 $\sum_{j\ge1}x^j/j=-\log(1-x)$。全部置换的生成函数为 $\exp(-\log(1-x))=1/(1-x)$，系数乘 $n!$ 后即 $n!$。错排禁止一元循环，所以生成函数为 $e^{-x}/(1-x)$，与式\eqref{eq:derangements}一致。这里指数与对数均用其形式级数定义，相应恒等式可由形式微分和常数项唯一性验证。

\begin{example}[单人和配对]
把 $n$ 个不同的人分成单人组或双人组，组件生成函数为 $x+x^2/2$，整体为 $\exp(x+x^2/2)$。若双人组数为 $j$，先把 $2j$ 人配对，再让其他人各自单独成组，数量为 $n!/((n-2j)!2^jj!)$。对 $0\le j\le\lfloor n/2\rfloor$ 求和；$n=4$ 时为 $1+6+3=10$。同样的数也计数只含一元、二元循环的置换。
\end{example}

\section{习题}
\begin{enumerate}
  \item\label{ex:egf-1} 用递推计算 $S(5,2)$、$S(5,3)$，再由满射计数核对前者。
  \item\label{ex:egf-2} 求 $B_5$，并解释递推中 $\binom nk$ 的标签选择含义。
  \item\label{ex:egf-3} 写出所有块大小均为偶数的集合分拆的指数生成函数，求 $n=4$ 时的数量。
  \item\label{ex:egf-4} 写出没有二元循环的置换的指数生成函数；保留一元循环，与错排条件作比较。
\end{enumerate}

\chapter{Catalan 数与整数分拆}
\label{ch:classical-structures}

本章研究两类重要对象。合法括号、受限路径和有序树共享递归分解，因此出现同一数列；整数分拆则忘掉部件顺序，适合用重数和图形表示。仅有前几项计数相同不足以证明两类对象等数，必须建立双射或共同的完整递推。

\section{合法括号、Dyck 路径与平面二叉树}

由 $n$ 个左括号和 $n$ 个右括号组成的字符串，若每个前缀中左括号不少于右括号，称为合法括号串。将左括号变成上升步 $(1,1)$，右括号变成下降步 $(1,-1)$，得到从 $(0,0)$ 到 $(2n,0)$ 且不低于横轴的\keyterm{Dyck 路径（Dyck path）}；按步骤还原括号给出逆映射。

非空合法括号串可唯一写成 $(A)B$，其中外层右括号是首个使左右括号数重新相等的位置，$A,B$ 各为合法括号串。把空串对应为一个叶子，非空串对应为一个根，其左右子树分别编码 $A,B$，便得到每个内部节点恰有两个有序孩子的满平面二叉树。大小 $n$ 对应内部节点数，左右次序属于结构，不能任意交换。

\section{Catalan 递推、生成函数与闭式}

令 $C_n$ 为含 $n$ 对括号的合法串数，$C_0=1$。上述唯一分解立即给出
\begin{equation}
 C_{n+1}=\sum_{i=0}^nC_iC_{n-i},\qquad
 C(x)=1+xC(x)^2,\qquad C(x)=\sum_{n\ge0}C_nx^n.
 \label{eq:catalan}
\end{equation}
前几项为 $1,1,2,5,14,42$。方程中常数项必须为 $1$，所以形式解为 $C(x)=(1-\sqrt{1-4x})/(2x)$，平方根取常数项为 $1$ 的唯一形式级数；分子常数项为零，故除以 $x$ 合法。另一个代数根不是普通幂级数。

\begin{theorem}[Catalan 公式]
对 $n\ge0$，$C_n=\frac1{n+1}\binom{2n}{n}$。
\end{theorem}
\begin{proof}
$n=0$ 时成立。对 $n\ge1$，不受高度限制且终点高度为零的路径有 $\binom{2n}{n}$ 条。对于曾到达高度 $-1$ 的坏路径，反转起点至首次到达 $-1$ 的全部上下步。这段原来下降步比上升步多一，反转后全路径有 $n+1$ 个上升步和 $n-1$ 个下降步，终点高度为 $2$。反之，任意终点为 $2$ 的路径必到达 $1$，反转其首次到达 $1$ 前的步骤便得到坏路径；两操作互逆。因此坏路径数为 $\binom{2n}{n-1}$。相减得
\[
 C_n=\binom{2n}{n}-\binom{2n}{n-1}
 =\frac1{n+1}\binom{2n}{n}.
\]
\end{proof}

这称为\keyterm{反射法（reflection principle）}。它的关键是指定首次触及边界的位置，保证反射可逆。凸 $(n+2)$ 边形的三角剖分也有 $C_n$ 种（$n\ge1$）：固定一条边，其所属三角形唯一地把剩余部分分成两个较小多边形；将空的边区域计为一种，所得递推与式\eqref{eq:catalan}相同。

\section{整数分拆与 Ferrers 图}

\begin{definition}
整数 $n\ge0$ 的\keyterm{分拆（partition）}是非增正整数序列 $\lambda_1\ge\cdots\ge\lambda_\ell>0$，满足 $\sum_i\lambda_i=n$；$n=0$ 时只有空分拆。总数记为 $p(n)$，其中 $p(0)=1$。各 $\lambda_i$ 称为部件。
\end{definition}

用左对齐的方格表示每个部件，第 $i$ 行有 $\lambda_i$ 格，所得图称为\keyterm{Ferrers 图（Ferrers diagram）}。交换行列得到共轭分拆 $\lambda'_j=|\{i:\lambda_i\ge j\}|$。转置两次恢复原图，且最大部件与部件个数交换，故“最大部件至多 $k$”与“部件个数至多 $k$”的分拆等数。

例如 $5$ 的分拆为 $5$、$4+1$、$3+2$、$3+1+1$、$2+2+1$、$2+1+1+1$、$1+1+1+1+1$，故 $p(5)=7$。整数合成保留部件顺序，共有 $16$ 个；集合分拆保留元素标签，计数又不同。三者不是同一公式的不同记号。

\section{分拆生成函数与基本恒等式}

部件 $j$ 的重数可以任取非负整数，对生成函数贡献 $1+x^j+x^{2j}+\cdots$，不同部件重数独立，因此
\begin{equation}
 \sum_{n\ge0}p(n)x^n=\prod_{j\ge1}\frac1{1-x^j},\qquad
 \prod_{j\ge1}(1+x^j)=\prod_{j\ge1}\frac1{1-x^{2j-1}}.
 \label{eq:partition-products}
\end{equation}
无穷乘积在形式意义下良定义，因为 $x^n$ 的系数只涉及 $j\le n$。第二式由 $1+x^j=(1-x^{2j})/(1-x^j)$ 逐系数约去偶数因子得到，说明不同部件分拆与只含奇数部件的分拆等数。

该结论还有保持总和的双射。对奇数部件 $u$ 的重数写二进制展开 $m_u=\sum_j\varepsilon_j2^j$，将其替换成所有 $\varepsilon_j=1$ 时的部件 $2^ju$。所得部件互不相同，因为每个正整数的奇数部分与二的指数唯一。逆向把每个部件 $2^ju$ 拆成 $2^j$ 个 $u$，即可恢复原分拆。例如 $3+3+1+1+1$ 对应 $6+2+1$。

\section{习题}
\begin{enumerate}
  \item\label{ex:classical-structures-1} 写出全部五个三对括号的合法串，给出其中一个对应的树和 Dyck 路径编码。
  \item\label{ex:classical-structures-2} 求凸七边形的三角剖分数，并说明参数为何不是 $C_7$。
  \item\label{ex:classical-structures-3} 求分拆 $(4,2,1)$ 的共轭，写出部件至多为 $3$ 的分拆生成函数。
  \item\label{ex:classical-structures-4} 用不同部件与奇数部件之间的双射，把 $8+6+3+1$ 转化为只含奇数部件的分拆，并核对总和。
\end{enumerate}

\chapter{图、路径与树的计数}
\label{ch:graphs-trees}

道路、通信线路和相互关系都能用图表达。图论关注顶点之间如何连接，组合计数则进一步询问满足指定结构的图有多少个。首先固定顶点与边的含义，再讨论连通、遍历和树的编码。

\section{图、度数与关联计数}

\begin{definition}
有限\keyterm{简单无向图（simple undirected graph）} $G=(V,E)$ 由有限顶点集 $V$ 及其二元素子集构成的边集 $E$ 组成。边 $\{u,v\}$ 也记为 $uv$；此时称 $u,v$ 邻接。顶点 $v$ 的邻点集记为 $N(v)$，度数记为 $d(v)=|N(v)|$。
\end{definition}

简单图没有环和重边。子图可删除顶点和边；由 $U\subseteq V$ 导出的诱导子图 $G[U]$ 则保留原图中两个端点都在 $U$ 的全部边。两个图之间若存在保持邻接与非邻接的顶点双射，称它们\keyterm{同构（isomorphic）}。固定顶点集 $[n]$ 时，每对顶点可独立决定是否连边，故有 $2^{\binom n2}$ 个标号简单图；按同构分类不能一般地直接除以 $n!$。

\begin{theorem}[握手定理]
$\sum_{v\in V}d(v)=2|E|$。因此奇度顶点的个数为偶数。
\end{theorem}
\begin{proof}
计数关联对 $(v,e)$，其中顶点 $v$ 是边 $e$ 的端点。按顶点计数得左侧，按边计数每边有两个端点，得右侧。度数总和为偶数，故其中奇数项的个数也为偶数。
\end{proof}

\section{连通性、路径与 Euler 道}

相邻顶点依次连接成\keyterm{游走（walk）}；边不重复的游走称为\keyterm{迹（trail）}，顶点不重复的游走称为\keyterm{路径（path）}。除首尾重合外顶点均不同的闭游走称为\keyterm{圈（cycle）}，在简单图中长度至少为 $3$。任意两点之间存在路径的非空图称为连通图，其极大连通部分称为连通分支。

经过每条边恰一次的迹称为\keyterm{Euler 道（Euler trail）}；首尾相同者为 Euler 闭迹。以下假定至少有一条边，且所有非孤立顶点处于同一连通分支。

\begin{theorem}
图存在 Euler 闭迹当且仅当所有顶点度数均为偶数；存在首尾不同的 Euler 道当且仅当恰有两个奇度顶点，且它们正是两个端点。
\end{theorem}
\begin{proof}
沿迹经过一个非端点时，每次进入都配对一次离开，所以其度数必须为偶数；首尾不同则两端各有一条未配对边，给出必要性。反之，若度数全偶，从有边顶点出发不断沿未用边前进，有限性保证最终无边可走。除起点外，进入顶点后已用边数在本次进入时为奇数，总度数为偶数，所以至少还有一条未用边，故只能在起点结束。若尚有未用边，由连通性，已有闭迹上必有顶点与未用边相接；删去已有闭迹后各点度数仍为偶数，从该点在剩余图中重复构造闭迹，再插入原闭迹。每次至少增加一条边，最终得到所需闭迹。若恰有两个奇度顶点，临时在它们之间增加一条允许重复的辅助边，构造闭迹后删除该边，即得开迹。
\end{proof}

\keyterm{Hamilton 圈（Hamilton cycle）}要求每个顶点恰经过一次，其任务与 Euler 道不同。两个三角形只共享一个顶点的图所有度数为偶数，因而有 Euler 闭迹，却没有 Hamilton 圈：若要访问两个三角形的其余顶点，必须重复经过公共顶点。度数奇偶条件不能作为 Hamilton 性的判据。

\section{树的等价刻画与生成树}

\begin{definition}
连通且无圈的图称为\keyterm{树（tree）}，无圈图称为\keyterm{森林（forest）}。度数为 $1$ 的顶点称为叶子。包含原图全部顶点的树子图称为\keyterm{生成树（spanning tree）}。
\end{definition}

\begin{theorem}
对含 $n\ge1$ 个顶点的简单图，下列条件等价：图是树；任意两点之间恰有一条路径；图连通且有 $n-1$ 条边。
\end{theorem}
\begin{proof}
连通保证路径存在。若两点间有两条不同路径，从首次分离到再次相遇的部分构成圈，因此树中路径唯一；反之，圈上两点有两条路径，故路径唯一蕴含无圈与连通。对 $n\ge2$ 的树，取最长路径，其端点不能有路径外邻点，否则可延长；也不能邻接路径上非相邻点，否则出现圈，故端点是叶子。删去叶子及其一条边后仍为树，归纳得边数为 $n-1$。最后，任意连通图可反复删去圈上的一条边而不破坏连通，终得生成树；若原图本来只有 $n-1$ 条边，则没有额外边可删，原图就是树。
\end{proof}

由每个分支为树，含 $n$ 个顶点、$c$ 个分支的森林有 $n-c$ 条边。一棵树选定一个顶点作根后，可按到根的距离定义父子关系；根树并不自动具有孩子之间的左右次序，只有再指定次序才是有序树。

\section{Pr\"ufer 编码与 Cayley 公式}

对顶点集 $[n]$、$n\ge2$ 的树，重复 $n-2$ 次以下操作：删除当前最小标签叶子，并记录它唯一邻点的标签。所得长度 $n-2$ 的串称为\keyterm{Pr\"ufer 序列（Pr\"ufer sequence）}。一个标签 $v$ 恰出现 $d(v)-1$ 次，因为删除其他叶子时，每记录一次 $v$ 就删掉一条与它相连的边，直到它自己成为叶子或留在最后两点中。

反向从任意串 $(a_1,\ldots,a_{n-2})\in[n]^{n-2}$ 出发。保留全部标签；第 $i$ 步取尚保留且未出现在后缀 $a_i,\ldots,a_{n-2}$ 中的最小标签 $b_i$，连边 $b_ia_i$ 并删除 $b_i$；最后连剩下两点。此时剩余标签数总比后缀长度多二，故 $b_i$ 必存在；$a_i$ 尚保留且不同于 $b_i$。逆序看这一构造，每步把一个新叶子接入已有树，所以最终为树。编码时，当前不再出现于后缀的标签恰是当前叶子，故两过程逐步互逆。

\begin{theorem}[Cayley 公式]
对 $n\ge2$，顶点集 $[n]$ 上的树有 $n^{n-2}$ 棵；$n=1$ 时有一棵。
\end{theorem}
\begin{proof}
编码是树与 $[n]^{n-2}$ 之间的双射，后者有 $n^{n-2}$ 个元素。单顶点图单独计数。
\end{proof}

例如序列 $(3,3)$ 对应 $[4]$ 上以 $3$ 为中心的星形树。进一步，若 $n\ge2$ 且各 $d_i\ge1$、$\sum_i d_i=2n-2$，指定度数序列的树数为 $(n-2)!/\prod_i(d_i-1)!$，因为标签 $i$ 须在序列中出现 $d_i-1$ 次。

\section{矩阵树定理与计数的线性代数表示}

设 $A$ 为邻接矩阵，$D=\operatorname{diag}(d(1),\ldots,d(n))$，$L=D-A$ 称为 Laplace 矩阵。记 $L^{(r)}$ 为删去第 $r$ 行及第 $r$ 列的矩阵。

\begin{theorem}[矩阵树定理（matrix-tree theorem）]
对 $n\ge2$ 的有限连通简单图，其生成树数为 $\det L^{(r)}$，与 $r$ 的选择无关。
\end{theorem}
\begin{proof}
任意给每条边一个方向，构造 $n\times |E|$ 的关联矩阵 $B$：每列在起点处为 $1$、终点处为 $-1$，其余为零。直接计算得 $BB^{\mathsf T}=L$。删去第 $r$ 行，得矩阵 $\widetilde B$，故 $L^{(r)}=\widetilde B\widetilde B^{\mathsf T}$。由附录\ref{app:prerequisites}中的 Cauchy--Binet 公式，
\[
 \det L^{(r)}=\sum_{\substack{S\subseteq E\\|S|=n-1}}
                 \det(\widetilde B_S)^2.
\]
若 $S$ 含圈，沿圈取带符号的列和为零，行列式为零。若无圈，由森林边数公式知 $(V,S)$ 连通，因而为树。此时取一个不是 $r$ 的叶子，沿其只有一个非零元的行展开，并递归删去叶子与对应边，最后得到 $|\det\widetilde B_S|=1$。所以每棵生成树贡献 $1$，其他边集贡献零。
\end{proof}

以上证明的行列式预备可与讲义\cite{matrix-tree-notes}对照。三角形的约化 Laplace 矩阵为 $\begin{pmatrix}2&-1\\-1&2\end{pmatrix}$，行列式为 $3$，与删去任意一边得到的三棵生成树一致。

\section{习题}
\begin{enumerate}
  \item\label{ex:graphs-trees-1} 证明树中至少有两个叶子（顶点数至少为 $2$），并说明最大度数不超过 $2$ 的树是什么图。
  \item\label{ex:graphs-trees-2} 解码顶点集 $[5]$ 上的 Pr\"ufer 序列 $(2,2,4)$，列出边集与各点度数。
  \item\label{ex:graphs-trees-3} 求 $[5]$ 上度数分别为 $(3,2,1,1,1)$ 的标号树数。
  \item\label{ex:graphs-trees-4} 用删除一条边的直接计数与矩阵树定理分别求四边形的生成树数。
\end{enumerate}

\chapter{匹配、染色与图的结构}
\label{ch:matching-coloring}

分配资源要求互不冲突地选择边，安排时间要求给相邻对象不同颜色。这些问题需要结构条件与构造方法共同作用。特别是，“再也不能直接加入一条边”只是一种局部停止条件，并不等于已经找到最多的边。

\section{二分图、匹配与 Hall 定理}

\begin{definition}
图 $G=(L\cup R,E)$ 若 $L,R$ 不交且每条边恰有一个端点在 $L$、一个在 $R$，称为\keyterm{二分图（bipartite graph）}。两两不共享端点的边集称为\keyterm{匹配（matching）}；每个 $L$ 中顶点都在某条匹配边上时，称该匹配饱和 $L$。对 $S\subseteq L$，记 $N(S)=\bigcup_{u\in S}N(u)$。
\end{definition}

\begin{theorem}[Hall 婚配定理]
有限二分图存在饱和 $L$ 的匹配，当且仅当对每个 $S\subseteq L$ 都有 $|N(S)|\ge|S|$。
\end{theorem}
\begin{proof}
必要性来自 $S$ 中各点必须匹配到不同邻点。充分性对 $|L|$ 归纳，$|L|=0,1$ 直接成立。若存在非空真子集 $S\subset L$ 满足 $|N(S)|=|S|$，则 $S$ 与 $N(S)$ 诱导的二分图满足 Hall 条件，由归纳可饱和 $S$。对剩余左侧点集的任意子集 $T$，由原图对 $S\cup T$ 的条件得
\[
 |N(T)\setminus N(S)|=|N(S\cup T)|-|N(S)|\ge|T|.
\]
故删去 $S\cup N(S)$ 的余图也满足条件，可饱和剩余左侧，两匹配合并即可。若不存在这样的 $S$，则每个非空真子集的邻点数至少多一。任选边 $uv$，其中 $u\in L$，删去 $u,v$ 后每个非空左侧子集至多失去一个邻点，仍满足 Hall 条件；由归纳取得匹配，再加回 $uv$。
\end{proof}

例如三名学生的候选题目集分别为 $\{a,b\},\{a\},\{b,c\}$，可取匹配 $1b,2a,3c$。若三人的候选集都是 $\{a,b\}$，每人虽都有选择，但取全部三人时 Hall 条件失败，故无法全部分配。条件必须对所有子集成立。

\section{增广路与最大匹配算法}

不能再添加一条边的匹配称为\keyterm{极大匹配（maximal matching）}；边数在全部匹配中最大者称为\keyterm{最大匹配（maximum matching）}。四点路径的中间一条边形成极大匹配，但两端的两条边构成更大匹配。

相对于匹配 $M$，一条边在 $M$ 内外交替、两个端点均未匹配的路径称为\keyterm{增广路（augmenting path）}。将其边的选中状态翻转后，仍为匹配，且大小增加一。反之，若另有 $|M'|>|M|$，考察对称差 $M\mathbin{\triangle}M'$；其中顶点度数至多为二，各分支为交替路径或偶圈。偶圈上两类边等多，而总数不等意味着某条路径的 $M'$ 边多一，该路径便是 $M$ 的增广路。因此匹配最大当且仅当不存在增广路。

在二分图中，从所有未匹配左侧顶点出发，沿未匹配边向右、匹配边向左搜索；到达未匹配右侧顶点便找到增广路。算法从空匹配开始，反复搜索并翻转，失败时停止。每轮增加一条匹配边，至多成功 $|V|/2$ 轮；邻接表实现的一轮搜索为 $O(|V|+|E|)$，故这一直接实现总成本为 $O(|V|(|V|+|E|))$。结论来自上述增广判据，不依赖遍历顺序。

\section{顶点覆盖与 K\H{o}nig 定理}

\keyterm{顶点覆盖（vertex cover）}是与每条边至少相交于一个端点的顶点集。每条匹配边都须由覆盖中的不同顶点照顾，故任意匹配大小不超过任意覆盖大小。

\begin{theorem}[K\H{o}nig 定理]
有限二分图的最大匹配大小等于最小顶点覆盖大小。
\end{theorem}
\begin{proof}
取最大匹配 $M$，按上一节的交替搜索规则，从所有未匹配左侧点出发，记可达集为 $Z$。令 $C=(L\setminus Z)\cup(R\cap Z)$。若某边 $uv$ 未被覆盖，则 $u\in L\cap Z$、$v\in R\setminus Z$；若该边不在 $M$ 中，可从 $u$ 到达 $v$，矛盾；若在 $M$ 中，已匹配的 $u$ 只能由其匹配伙伴 $v$ 到达，也矛盾。故 $C$ 为覆盖。

不存在可达的未匹配右侧点，否则得到增广路；所有未匹配左侧点都在 $Z$。每条匹配边的两个端点要么都可达，要么都不可达，因而恰有一个端点在 $C$ 中，且 $C$ 不含未匹配点。于是 $|C|=|M|$，结合一般不等式得到结论。
\end{proof}

同时给出大小相同的匹配与覆盖，就是最优性证书。二分图假设不可删除：三角形最大匹配大小为 $1$，最小顶点覆盖大小为 $2$。

\section{图染色与染色多项式}

将顶点赋予 $q$ 种有名称颜色，使邻接点异色，称为\keyterm{正常顶点染色（proper vertex coloring）}。最少所需颜色数称为色数 $\chi(G)$。若最大度数为 $\Delta$，依次给各点选一种尚未被已染邻点使用的颜色，则 $\Delta+1$ 种颜色总够用，因为被禁用的颜色至多为 $\Delta$ 种。这给出上界及构造，但一般不保证最少用色。

记 $P_G(q)$ 为使用 $q$ 种有名称颜色的正常染色数。若 $e=uv$ 是一条边，$G-e$ 的正常染色按 $u,v$ 是否同色分类：异色者对应 $G$ 的染色，同色者对应把 $u,v$ 合并所得 $G/e$ 的染色。缩边产生的平行边可合并，所以
\[
 P_G(q)=P_{G-e}(q)-P_{G/e}(q).
\]
从无边图的 $q^{|V|}$ 出发归纳，证明这是关于 $q$ 的多项式，称为\keyterm{染色多项式（chromatic polynomial）}。$n\ge1$ 的树先染根，再沿树为每个孩子避开父色，得 $q(q-1)^{n-1}$；$n\ge3$ 的圈由删边--缩边归纳得 $(q-1)^n+(-1)^n(q-1)$。计数不对颜色名称置换取商。

\section{平面图与 Euler 公式入门}

能在平面中画出且边仅在共同端点相交的图称为\keyterm{可平面图（planar graph）}；固定这样一个嵌入后称为平面图（plane graph）。平面被嵌入分成若干面，包括无界外部面。以下使用简单闭曲线将平面分成内外两部分的拓扑事实；这是 Jordan 曲线定理的内容，此处作为已知事实，不在组合教材中证明。

对连通平面图，若 $n,m,f$ 分别为顶点、边、面数，则
\begin{equation}
 n-m+f=2.
 \label{eq:planar-euler}
\end{equation}
证明可不断删去圈上的边：连通性不变，边两侧不同的面合并，$m,f$ 各减一，故左侧不变；最终为生成树，其 $m=n-1,f=1$，得到等式。这里关于圈边分隔面的判断正是所用平面拓扑事实。

对 $n\ge3$ 的连通简单平面图，按面边界计数时桥在同一面上计两次，每个面的边界长度至少为三，因此 $3f\le2m$。代入式\eqref{eq:planar-euler}得 $m\le3n-6$。若还是二分图，则没有奇圈，各面边界长度至少为四，得 $m\le2n-4$。不连通情形可先在各分支之间添加不相交的边使其连通，保留相应性质后应用上界。$K_5$ 有 $10>9$ 条边，$K_{3,3}$ 有 $9>8$ 条边，故均不可平面；边数满足这些上界则仅是必要条件。

\section{习题}
\begin{enumerate}
  \item\label{ex:matching-coloring-1} 三名学生的候选集合为 $\{a,b\},\{a,b\},\{b,c\}$，给出饱和匹配，并验证 Hall 条件。
  \item\label{ex:matching-coloring-2} 在四点路径上先选中间边，写出一条增广路和翻转后的匹配。
  \item\label{ex:matching-coloring-3} 计算四边形的染色多项式及使用三种有名称颜色的正常染色数。
  \item\label{ex:matching-coloring-4} 证明任意非空简单可平面图有一个度数不超过 $5$ 的顶点，并由删除顶点归纳构造六色染色。
\end{enumerate}

\chapter{群作用与对称计数}
\label{ch:symmetry}

四个位置的二色染色有 $2^4=16$ 种，但把它们排在正方形上并识别旋转后只剩六种。不能用 $16/4$ 计算，因为单色图案旋转任意角度都不变，而其他图案的旋转副本数可能为二或四。群作用提供了处理不同重复次数的统一方法。

\section{对称变换、群与轨道}

\begin{definition}
集合 $G$ 配备满足结合律的乘法，具有单位元 $e$ 且每个元素都有逆元，称为\keyterm{群（group）}。群 $G$ 在集合 $X$ 上的\keyterm{作用（action）}是映射 $(g,x)\mapsto g\cdot x$，满足 $e\cdot x=x$ 与 $(gh)\cdot x=g\cdot(h\cdot x)$。
\end{definition}

每个群元素都诱导 $X$ 的一个置换，因为 $g^{-1}$ 给出其逆映射。正 $n$ 边形的旋转形成循环群，$n\ge3$ 时旋转与反射合起来形成含 $2n$ 个元素的二面体群。这里变换作用于位置，再由位置置换诱导图案变换；颜色名称保持不变。

对 $x\in X$，\keyterm{轨道（orbit）}为 $G\cdot x=\{g\cdot x:g\in G\}$，\keyterm{稳定子（stabilizer）}为 $G_x=\{g:g\cdot x=x\}$；对 $g\in G$，固定点集为 $\operatorname{Fix}(g)=\{x:g\cdot x=x\}$。处于同一轨道是等价关系，轨道集合记为 $X/G$，它正是识别指定对称后的方案集合。稳定子固定一个对象，固定点集固定一个变换，二者不可混淆。

\section{轨道--稳定子关系与平均固定点}

\begin{proposition}
若 $G$ 有限，则 $|G\cdot x|=|G|/|G_x|$。
\end{proposition}
\begin{proof}
考虑满射 $g\mapsto g\cdot x$。若 $g_0\cdot x=y$，则 $g\cdot x=y$ 当且仅当 $g_0^{-1}g\in G_x$，所以该纤维恰为 $g_0G_x$，大小为 $|G_x|$。应用等重计数原理即得。
\end{proof}

\begin{theorem}[Burnside 引理]
有限群 $G$ 作用于有限集合 $X$ 时，
\begin{equation}
 |X/G|=\frac1{|G|}\sum_{g\in G}|\operatorname{Fix}(g)|.
 \label{eq:burnside}
\end{equation}
\end{theorem}
\begin{proof}
对关联对集合 $\{(g,x):g\cdot x=x\}$，先按 $g$ 求和得到右侧分子，按 $x$ 求和得到 $\sum_x|G_x|$。同一轨道 $O$ 中每个点的稳定子大小为 $|G|/|O|$，所以整条轨道贡献 $|G|$。轨道总贡献为 $|G||X/G|$，相除即得。
\end{proof}

所有非单位元都不固定任何对象时，轨道均有 $|G|$ 个元素，公式才退化为简单除以群大小。一般情形平均的是各变换的固定点数，并非任意猜测一种代表的重复次数。

\section{项链、手链与置换循环}

设 $g$ 是位置集上的置换，共有 $c(g)$ 个循环。一个用 $q\ge1$ 种有名称颜色染色的图案被 $g$ 固定，当且仅当每个循环上的颜色恒定；各循环独立选色，故固定图案数为 $q^{c(g)}$。

在 $n$ 个圆周位置中，转动 $j$ 格产生 $\gcd(n,j)$ 个循环，约定 $\gcd(n,0)=n$。只识别旋转的\keyterm{项链（necklace）}数为
\[
 \frac1n\sum_{j=0}^{n-1}q^{\gcd(n,j)}.
\]
再识别反射得到\keyterm{手链（bracelet）}。当 $n\ge3$ 为奇数时，每个反射有一个不动位置和 $(n-1)/2$ 对位置；当 $n$ 为偶数时，$n/2$ 个反射固定两个位置，其余 $n/2$ 个不固定位置。因此在旋转固定点总和之外，反射部分分别贡献 $nq^{(n+1)/2}$ 或 $\frac n2(q^{n/2+1}+q^{n/2})$，总和除以 $2n$ 即得。

\begin{example}[正方形顶点染色]
正方形的旋转固定点总和为 $q^4+q^2+2q$，故旋转等价类数为 $(q^4+q^2+2q)/4$。四个反射中两个固定一对对角顶点，贡献 $2q^3$；两个交换两对顶点，贡献 $2q^2$，故完整对称下有 $(q^4+2q^3+3q^2+2q)/8$ 类。$q=2$ 时两者都为 $6$；$q=3$ 时分别为 $24$ 与 $21$。两种等价关系在一个小例中数值相同，并不意味着它们总相同。
\end{example}

\section{循环指标与 P\'olya 计数}

设 $G$ 作用于 $n$ 个位置，$c_i(g)$ 为 $g$ 的长度为 $i$ 的循环数。定义\keyterm{循环指标（cycle index）}
\begin{equation}
 Z_G(s_1,\ldots,s_n)=\frac1{|G|}\sum_{g\in G}\prod_{i=1}^n s_i^{c_i(g)}.
 \label{eq:cycle-index}
\end{equation}
将所有 $s_i$ 代为 $q$ 就得到不限颜色库存的轨道数。若用变量 $y_a$ 标记颜色 $a$ 的使用次数，一个长度为 $i$ 的循环被固定时必须整体染成某种颜色，贡献 $y_1^i+\cdots+y_q^i$。

\begin{theorem}[P\'olya 计数的库存形式]
颜色名称固定时，每个染色轨道按其颜色使用次数贡献一个单项式，所得多项式为
\[
 Z_G\left(\sum_{a=1}^q y_a,\sum_{a=1}^q y_a^2,\ldots,
                 \sum_{a=1}^q y_a^n\right).
\]
因此指定颜色库存的轨道数为相应单项式的系数。
\end{theorem}
\begin{proof}
对每个 $g$，固定染色的权重和为上述循环因子的乘积。权重在同一轨道上不变，故在 Burnside 双重计数中，一条轨道连同其稳定子贡献 $|G|$ 倍该轨道的权重。对 $g$ 的权重和取平均，便使每条轨道恰贡献一次。
\end{proof}

正方形旋转群的循环指标为 $(s_1^4+s_2^2+2s_4)/4$。两色库存各为二时，$[y_1^2y_2^2]$ 的值为 $(6+2+0)/4=2$，对应同色位置相邻和同色位置相对。限制库存后，不能仍以 $q^{c(g)}$ 作为固定点数。

\section{习题}
\begin{enumerate}
  \item\label{ex:symmetry-1} 求五个圆周位置的二色项链数，再加入反射等价，比较结果。
  \item\label{ex:symmetry-2} 写出正方形二面体群的循环指标，并求三色染色轨道数。
  \item\label{ex:symmetry-3} 求六个圆周位置上恰有两个红色、四个蓝色的旋转等价类数。
  \item\label{ex:symmetry-4} 给出一个群作用使不同轨道大小不同，逐个列出稳定子并验证 Burnside 引理。
\end{enumerate}

\chapter{偏序集与 M\"obius 反演}
\label{ch:posets}

子集包含和整数整除都允许比较某些对象，却不要求任意两个对象可比。偏序描述这种结构，而 M\"obius 反演从一般的偏序求和中恢复原始数据。容斥只是这一机制在子集偏序上的一个实例。

\section{偏序、链、反链与 Hasse 图}

\begin{definition}
集合 $P$ 上的关系 $\le$ 若具有自反性、反对称性和传递性，称为\keyterm{偏序（partial order）}，$(P,\le)$ 称为偏序集。任意两元素可比的子集称为\keyterm{链（chain）}；任意两个不同元素均不可比的子集称为\keyterm{反链（antichain）}。
\end{definition}

极小元是没有严格更小元素的元素；最小元则小于等于每个元素。极小元可有多个，最小元若存在则唯一。若 $x<y$ 且不存在 $x<z<y$，称 $y$ 覆盖 $x$；Hasse 图只画覆盖关系，并把较大的元素画在上方，通过向上的路径恢复整个偏序。

$2^{[n]}$ 配备包含关系称为布尔格。若偏序集中每对元素都有最大下界与最小上界，称为\keyterm{格（lattice）}；在布尔格中它们分别为交集与并集。正整数 $N$ 的因子集按整除排序也构成格，最大下界为最大公因数，最小上界为最小公倍数，二者仍是 $N$ 的因子。

\section{关联代数与 M\"obius 函数}

本章只处理有限偏序集。对定义在可比对 $x\le y$ 上的实值函数，定义卷积
\[
 (f*g)(x,y)=\sum_{x\le z\le y}f(x,z)g(z,y).
\]
交换有限求和可验证结合律；单位元为 $\delta(x,y)$，当 $x=y$ 时为 $1$，否则为零。令 $\zeta(x,y)=1$，并递归定义
\[
 \mu(x,x)=1,\qquad
 \mu(x,y)=-\sum_{x\le z<y}\mu(x,z)\quad(x<y).
\]
递归按区间元素数进行，每次用到的区间严格较小，所以定义唯一。它直接给出 $\mu*\zeta=\delta$。

不断删除一个极小元，可以把有限偏序集列成与偏序相容的线性顺序。按此顺序把上述函数写成上三角矩阵，卷积对应矩阵乘法；$\zeta$ 的对角元全为 $1$，故可逆。于是右逆也是左逆，得到 $\zeta*\mu=\delta$。这类函数与卷积形成\keyterm{关联代数（incidence algebra）}，$\mu$ 称为该偏序集的 M\"obius 函数。

\section{反演公式与容斥的统一解释}

\begin{theorem}[M\"obius 反演]
有限偏序集上，若 $F(x)=\sum_{y\le x}f(y)$，则
\[
 f(x)=\sum_{y\le x}\mu(y,x)F(y).
\]
\end{theorem}
\begin{proof}
右侧代入 $F$ 并换序，得 $\sum_{z\le x}f(z)\sum_{z\le y\le x}\mu(y,x)$。内层和为 $(\zeta*\mu)(z,x)=\delta(z,x)$，故仅保留 $f(x)$。
\end{proof}

在布尔格上，$\mu(A,B)=(-1)^{|B\setminus A|}$：固定差集大小后，递归所需的和就是二项式交错和。因而
\[
 F(S)=\sum_{T\subseteq S}f(T)
 \quad\Longrightarrow\quad
 f(S)=\sum_{T\subseteq S}(-1)^{|S\setminus T|}F(T).
\]
按某种性质的实际支撑集分类，便回到容斥；若计数只依赖支撑大小，则再按大小汇总得到二项式反演。

在整数整除偏序上，令数论 M\"obius 函数 $\mu_{\mathrm{ar}}(m)$ 在 $m$ 含平方素因子时为零，在 $m$ 是 $r$ 个不同素数之积时为 $(-1)^r$，并置 $\mu_{\mathrm{ar}}(1)=1$。利用素因数唯一分解，$\sum_{d\mid m}\mu_{\mathrm{ar}}(d)$ 在 $m>1$ 时等于 $(1-1)^r=0$，在 $m=1$ 时为 $1$，故偏序函数满足 $\mu(a,b)=\mu_{\mathrm{ar}}(b/a)$。例如在 $12$ 的因子偏序中，$\mu(1,12)=0$，而 $\mu(2,12)=1$。

\section{Sperner 定理、LYM 不等式与链分解}

\begin{theorem}
设 $\mathcal A$ 为 $2^{[n]}$ 中的反链，则
\begin{equation}
 \sum_{A\in\mathcal A}\binom n{|A|}^{-1}\le1,
 \qquad |\mathcal A|\le\binom n{\lfloor n/2\rfloor}.
 \label{eq:sperner}
\end{equation}
第一式称为 LYM 不等式，第二式为 Sperner 定理。
\end{theorem}
\begin{proof}
每个 $[n]$ 的排列按前缀生成一条从空集到全集的最大链，反链至多与它相交一次。固定 $A$ 出现在此链中，当且仅当前 $|A|$ 个位置恰是 $A$ 的元素，共有 $|A|!(n-|A|)!$ 个排列。对所有排列与其链上的反链元素计数，得 $\sum_{A\in\mathcal A}|A|!(n-|A|)!\le n!$，相除即得第一式。二项式系数在中间层最大，故每个倒数至少为 $1/\binom n{\lfloor n/2\rfloor}$，推出第二式。全体大小为 $\lfloor n/2\rfloor$ 的子集构成达到上界的反链。
\end{proof}

另一个结构结论是 Dilworth 定理：有限偏序集的最大反链大小等于把全部元素划分成链所需的最少条数。给出一个基于匹配的证明。将 $n=|P|$ 个元素复制为左右两份，当 $x<y$ 时连边 $x_Ly_R$；设最大匹配大小为 $m$。匹配在原偏序中形成每点入度、出度至多为一的有向边，而且不可能有圈，故分成 $n-m$ 条覆盖全部元素的链。另一方面，由 K\H{o}nig 定理有大小为 $m$ 的顶点覆盖；取左右两个副本都不在覆盖中的原元素，至少有 $n-m$ 个，且构成反链，否则对应边未被覆盖。任意反链与每条链至多相交一次，因而两界相合，定理得证。

\section{习题}
\begin{enumerate}
  \item\label{ex:posets-1} 在 $\{2,3,6\}$ 的整除偏序中，找出极小元、最小元、极大元与最大元。
  \item\label{ex:posets-2} 列出 $6$ 的全部正因子，计算各区间的 M\"obius 值，并验证反演公式的一个实例。
  \item\label{ex:posets-3} 在 $2^{[5]}$ 中构造最大反链，并说明最大大小。
  \item\label{ex:posets-4} 将 $\{1,2,3,4,6,12\}$ 按整除关系划分成最少条链，给出同样大小的反链作为最优性证书。
\end{enumerate}

\chapter{极值组合与 Ramsey 理论}
\label{ch:extremal-ramsey}

精确计数研究合法对象总共有多少；极值组合则研究在避免某种结构的条件下，对象最多可以多大。Ramsey 理论从另一方向说明：规模足够大时，某些有序子结构无法避免。两种问题都需要严格区分构造给出的下界与普遍论证给出的上界。

\section{极值函数与上下界的分工}

\begin{definition}
对固定简单图 $H$，记 $\operatorname{ex}(n,H)$ 为不含 $H$ 同构子图的 $n$ 顶点简单图的最大边数，称为相应的\keyterm{极值数（extremal number）}。此处不要求被禁止的子图为诱导子图。
\end{definition}

给出一个有 $m$ 条边的 $H$-无关图，只能推出 $\operatorname{ex}(n,H)\ge m$；证明每个此类图至多有 $m$ 条边，才能推出反向不等式。两者结合才确定精确值。有限性保证最大值存在，因为固定标号集上的简单图只有有限个。

类似地，可研究 $[n]$ 的子集族在不出现某种包含关系时最多有多少成员；Sperner 定理就是禁止两个成员严格包含时的极值结论。不同问题的共同模式是“限制结构---构造候选---证明上界---判断能否相合”。

\section{Mantel 定理与 Tur\'an 图}

\begin{theorem}[Mantel 定理]
对 $n\ge1$，$\operatorname{ex}(n,K_3)=\lfloor n^2/4\rfloor$，其中 $K_r$ 表示 $r$ 个顶点的完全图。
\end{theorem}
\begin{proof}
若无三角形，则每条边 $uv$ 的两个邻点集不交，所以 $d(u)+d(v)\le n$。记边数为 $m$，对边求和得 $\sum_vd(v)^2\le mn$。另一方面，$\sum_v(d(v)-2m/n)^2\ge0$ 给出 $\sum_vd(v)^2\ge4m^2/n$。若 $m>0$，相合得 $m\le n^2/4$；$m=0$ 时同样成立。两部分大小为 $\lfloor n/2\rfloor,\lceil n/2\rceil$ 的完全二分图无三角形，且有 $\lfloor n^2/4\rfloor$ 条边，达到上界。
\end{proof}

把 $n$ 个顶点分成 $r\ge1$ 部，部内无边、部间全连，得到完全 $r$ 部图。若部大小为 $n_1,\ldots,n_r$，允许空部，则边数为 $(n^2-\sum_i n_i^2)/2$。当两个部大小相差至少二时，从大部移一个顶点到小部使平方和减小，故边数最大时各部大小相差至多一。此平衡图记为 $T_r(n)$，边数记为 $t_r(n)$。写 $n=qr+s$、$0\le s<r$，有
\[
 t_r(n)=\frac12\bigl(n^2-s(q+1)^2-(r-s)q^2\bigr).
\]

\begin{theorem}[Tur\'an 定理]
对 $n\ge0,r\ge1$，$\operatorname{ex}(n,K_{r+1})=t_r(n)$。
\end{theorem}
\begin{proof}
$T_r(n)$ 中任意团从每部至多取一点，故不含 $K_{r+1}$，给出下界。对 $r$ 归纳，并在固定 $r$ 时对 $n$ 归纳证明上界。$r=1$ 时图无边；$n<r$ 时 $t_r(n)=\binom n2$，结论成立。设 $r\ge2,n\ge r$。若图不含 $K_r$，由较小 $r$ 的结论，边数至多 $t_{r-1}(n)\le t_r(n)$，后一个不等式来自完全多部图的平衡极值。若含一个 $r$ 点团，则每个外部顶点至多与其中 $r-1$ 点邻接，否则出现 $K_{r+1}$。删除该团后，由较小 $n$ 的归纳假设，
\[
 |E|\le t_r(n-r)+(r-1)(n-r)+\binom r2=t_r(n).
\]
最后等号可通过向 $T_r(n-r)$ 的每一部各加一个顶点直接计算；平衡性保持不变。这完成双重归纳。
\end{proof}

定理给出极值数及达到它的构造；关于所有等号图的分类需进一步论证，本节不由边数公式直接声称这种分类。更多禁子图问题可参考\cite{zhao-graph-notes}。

\section{Ramsey 数与有限性证明}

\begin{definition}
对 $r,s\ge2$，\keyterm{Ramsey 数（Ramsey number）} $R(r,s)$ 是使得每个 $K_n$ 的红蓝边染色都含红色 $K_r$ 或蓝色 $K_s$ 的最小 $n$。其存在性由下面的上界证明。
\end{definition}

边界为 $R(2,s)=s$、$R(r,2)=r$：若有红边即出现红色 $K_2$，否则全部边蓝色；反向下界取少一个顶点的全蓝图即可。对 $r,s\ge3$，有
\begin{equation}
 R(r,s)\le R(r-1,s)+R(r,s-1)
       \le\binom{r+s-2}{r-1}.
 \label{eq:ramsey-recursion}
\end{equation}
证明第一步时，令 $a=R(r-1,s)$、$b=R(r,s-1)$，在 $K_{a+b}$ 中取顶点 $v$。其 $a+b-1$ 个邻点中，或者红邻点至少为 $a$，或者蓝邻点至少为 $b$。前一种情形在红邻点中得到蓝色 $K_s$ 或红色 $K_{r-1}$，后者连同 $v$ 构成红色 $K_r$；另一种情形对称。由 $r+s$ 归纳可同时证明这些数存在，第二步再用 Pascal 恒等式及边界归纳。

把红边看成某个图的边，蓝边就是补图的边，所以同一结论可表述为：足够大的图含 $r$ 点团或 $s$ 点独立集。\keyterm{独立集（independent set）}是内部没有边的顶点集，与团分别代表稠密和稀疏的局部结构。

\section{小规模精确值与一般界}

\begin{example}
$R(3,3)=6$。在六点完全图中，某顶点至少有三个同色邻点，设为红邻点。如果三点之间有红边，便与原顶点组成红三角形；否则三点自身组成蓝三角形，故 $R(3,3)\le6$。在五点上把五边形的边染红、其余边染蓝，两种颜色都构成五圈，均无三角形，故 $R(3,3)>5$。
\end{example}

对更大参数，构造一个避免两种单色团的染色给出下界，递推或其他结构定理给出上界。不能把上界当作精确值，也不能因为一种搜索未找到反例就认为上界已经达到。下一章的随机染色将提供无需直接展示图案的下界证明；这仍是严格的存在性论证，但通常不是具体构造。

\section{习题}
\begin{enumerate}
  \item\label{ex:extremal-ramsey-1} 七个顶点的无三角形图最多有多少条边？给出达到上界的图。
  \item\label{ex:extremal-ramsey-2} 计算 $t_3(8)$，明确三个部分的大小，并验证它不含 $K_4$。
  \item\label{ex:extremal-ramsey-3} 仅用 $R(3,3)=6$ 与式\eqref{eq:ramsey-recursion}求出 $R(3,4)$ 的一个上界；说明为何不能据此声称等号。
  \item\label{ex:extremal-ramsey-4} 把“六人中必有三人两两相识或三人两两不相识”翻译成图论命题，列出对相识关系所作的假设。
\end{enumerate}

\chapter{概率方法}
\label{ch:probabilistic}

概率方法通过随机选择证明确定性对象存在。若坏对象出现的期望数量小于一，就必有一次选择完全没有坏对象；若一个得分的期望很高，就必有某次选择达到至少该得分。随机性来自我们设计的有限样本空间，而非未知的自然分布。

\section{有限概率空间与指示变量}

设有限非空集合 $\Omega$ 上给定概率 $p(\omega)\ge0$，且 $\sum_{\omega\in\Omega}p(\omega)=1$。事件 $A\subseteq\Omega$ 的概率为 $\mathbb P(A)=\sum_{\omega\in A}p(\omega)$；实值\keyterm{随机变量（random variable）} $X$ 是 $\Omega$ 上的函数，其\keyterm{期望（expectation）}为 $\mathbb E X=\sum_\omega X(\omega)p(\omega)$。事件的指示变量 $\mathbf1_A$ 在 $A$ 中取一，否则取零。

由定义直接得到 $\mathbb E\mathbf1_A=\mathbb P(A)$。对有限个随机变量 $X_1,\ldots,X_m$ 和实常数 $a_i$，交换有限求和即得
\[
 \mathbb E\left(\sum_{i=1}^m a_iX_i\right)=\sum_{i=1}^m a_i\mathbb E X_i.
\]
这不要求独立。若事件 $A,B$ 独立，才有 $\mathbb P(A\cap B)=\mathbb P(A)\mathbb P(B)$；涉及多个事件的乘积时须有相应的联合独立性，不能由两两独立自动替代。

\section{第一矩方法与存在性}

\begin{proposition}[第一矩方法（first moment method）]
若 $X$ 为非负整数值随机变量且 $\mathbb E X<1$，则 $\mathbb P(X=0)>0$。任意有限实值随机变量还至少有一个正概率取值不小于 $\mathbb E X$。
\end{proposition}
\begin{proof}
由 $X\ge\mathbf1_{\{X\ge1\}}$，得 $\mathbb P(X\ge1)\le\mathbb E X<1$，故第一项成立。若所有正概率取值都严格小于期望，它们的加权平均也严格小于期望，与定义矛盾。
\end{proof}

更一般地，对非负 $X$ 和 $a>0$，$a\mathbf1_{\{X\ge a\}}\le X$ 给出 Markov 不等式 $\mathbb P(X\ge a)\le\mathbb E X/a$。事件并的指示变量至多为各事件指示变量之和，因此 $\mathbb P(\bigcup_iA_i)\le\sum_i\mathbb P(A_i)$。这些上界即使在事件相关时仍成立。

\begin{example}[大割的存在]
图有 $m$ 条边，将每个顶点独立等概率放入左右两组。对每条边 $e$，令 $I_e$ 表示其端点是否被分开，则 $\mathbb E I_e=1/2$。跨组边数 $X=\sum_eI_e$ 的期望为 $m/2$，故存在至少含 $\lceil m/2\rceil$ 条跨组边的划分，称为一个\keyterm{割（cut）}。证明使用期望线性性，不需要不同边的指示变量独立。
\end{example}

\section{随机染色与 Ramsey 下界}

将 $K_n$ 的每条边独立等概率染成红或蓝，设 $n\ge k\ge2$，令 $X$ 为单色 $K_k$ 的数量。每个 $k$ 元顶点集有 $\binom k2$ 条边，全红和全蓝是两个互斥事件，各自概率为 $2^{-\binom k2}$。因此
\begin{equation}
 \mathbb E X=\binom nk\,2^{1-\binom k2}.
 \label{eq:ramsey-first-moment}
\end{equation}
若右侧小于一，则存在没有单色 $K_k$ 的染色，所以 $R(k,k)>n$。例如 $k=5,n=10$ 时，期望为 $252/512<1$，得 $R(5,5)>10$。这个例子的意义是展示方法，而不是声称得到最佳下界。

不同顶点集上的单色团事件通常相关，但求和只需各自的概率。必须要求期望严格小于一；若等于一，上述论证无法排除 $X$ 恒为一的情况。

\section{删改法与条件期望法}

\keyterm{删改法（alteration method）}先随机产生候选结构，再删除局部障碍。设图有 $n$ 个顶点、$m$ 条边，各点独立以概率 $p\in[0,1]$ 入选。令 $X$ 为入选点数、$Y$ 为两端都入选的边数，则 $\mathbb E(X-Y)=np-mp^2$。对每条仍存在的入选边删去一个端点，至多删去初始的 $Y$ 个点，得到大小至少为 $X-Y$ 的独立集。因某次取值不低于期望，必存在大小至少为 $np-mp^2$ 的独立集；大小为整数时可向上取整。例如 $n=100,m=200,p=1/4$ 给出至少 $13$ 个点的独立集。删改保证是确定性的，随机性只用于选取较好的初始样本。

\keyterm{条件期望法（method of conditional expectations）}可以把某些存在性证明改为构造。对已固定的部分选择，剩余随机选择的平均得分是当前条件期望；若下一步有若干结果，它们的条件期望按概率平均等于当前值，因此至少有一个结果不降低它。逐步选择这样的结果，最终确定的得分不低于最初期望。

在大割问题中，按固定次序处理顶点，把当前顶点放到使它与已放置邻点之间的跨组边最多的一侧；与未放置点相关的边仍有一半的预期贡献，不受本次选择影响。于是条件期望不下降，最后割大小至少 $m/2$。邻接表中每条边只需常数次检查，可在 $O(n+m)$ 时间完成。一般问题只有在条件期望可有效计算时，此方法才给出高效算法。

\section{习题}
\begin{enumerate}
  \item\label{ex:probabilistic-1} 在均匀随机排列中，令 $X$ 为固定点数，证明 $n\ge1$ 时 $\mathbb E X=1$。这能否仅靠第一矩方法证明错排存在？
  \item\label{ex:probabilistic-2} 用随机二染色证明：若一个超图有 $m$ 条边、每条边恰含 $k\ge1$ 个顶点，且 $m2^{1-k}<1$，则存在无单色边的顶点染色。超图的边在此指顶点的子集。
  \item\label{ex:probabilistic-3} 对固定 $n,m>0$，在 $0\le p\le1$ 上最大化删改法下界 $np-mp^2$，分别讨论 $2m\ge n$ 与 $2m<n$。
  \item\label{ex:probabilistic-4} 在五边形上执行条件期望的大割算法，给出顶点处理次序、每步选择和最终割，说明它不必预先知道最优值。
\end{enumerate}

\chapter{组合设计与纠错编码}
\label{ch:designs-codes}

组合设计要求局部关联均匀，例如每对试验对象共同出现相同次数；纠错编码则要求信息表示彼此足够分离，以便容忍传输错误。前者通过双重计数约束参数，后者通过距离与不相交球约束规模。

\section{区组设计与参数约束}

\begin{definition}
设 $V$ 为 $v$ 点集合，$2\le k\le v$、$\lambda\ge1$。一族不同的 $k$ 元子集称为一个 $2$-$(v,k,\lambda)$ \keyterm{区组设计（block design）}，若每对不同点恰出现在 $\lambda$ 个子集中。各子集称为区组，其数量记为 $b$。
\end{definition}

固定一个点 $x$，计数包含 $x$ 的区组与其中另一个点形成的关联对。若 $x$ 出现在 $r_x$ 个区组中，该数为 $r_x(k-1)$；先选另一点，该数为 $\lambda(v-1)$。所以 $r_x$ 与 $x$ 无关，统一记为 $r$。再计数点与区组的关联对，得到
\begin{equation}
 r(k-1)=\lambda(v-1),\qquad vr=bk.
 \label{eq:design-parameters}
\end{equation}
由此 $r,b$ 必须是整数，并且在区组不重复的约定下 $b\le\binom vk$。这些是必要条件，仅有参数通过检查并不能替代实际构造。

还有一个有用的限制：若 $k<v$，则 $b\ge v$，称为 Fisher 不等式。令 $M$ 为 $v\times b$ 的零一关联矩阵，每列记录一个区组，则 $MM^{\mathsf T}=(r-\lambda)I+\lambda J$，其中 $J$ 为全一矩阵。由式\eqref{eq:design-parameters}及 $k<v$ 得 $r>\lambda$；任意非零实向量 $z$ 满足
\[
 z^{\mathsf T}MM^{\mathsf T}z=(r-\lambda)\sum_i z_i^2+
                       \lambda\left(\sum_i z_i\right)^2>0.
\]
故该矩阵秩为 $v$，而秩至多为 $b$，证明不等式。

\section{小设计与有限几何的例子}

七点集合上的七个区组
\[
 123,\quad145,\quad167,\quad246,\quad257,\quad347,\quad356
\]
构成 $2$-$(7,3,1)$ 设计，这里 $abc$ 简写集合 $\{a,b,c\}$。每个区组贡献三对点，七个区组共贡献 $21=\binom72$ 对，逐项可核查这些点对互不重复；因此每对恰出现一次。由式\eqref{eq:design-parameters}得 $r=3,b=7$，与逐点核对相符。

这一结构称为\keyterm{Fano 平面（Fano plane）}。令 $\mathbb F_2=\{0,1\}$，加法与乘法均模二进行；更系统的构造把七点视为 $\mathbb F_2^3$ 的七个非零向量，向量运算逐坐标进行，每条线为 $\{x,y,x+y\}$，其中 $x\ne y$。任意两不同非零向量确定唯一这样的三元组，所以关联条件自动成立。取编号 $1=100,2=010,3=110,4=001,5=101,6=011,7=111$，正好得到上面的区组。有限几何中的“线”是关联对象，不要求在普通平面上都画成直线。

\section{拉丁方与正交性}

\begin{definition}
$n\times n$ 表格使用 $n$ 个符号，每行每列各恰出现一次每个符号，称为 $n$ 阶\keyterm{拉丁方（Latin square）}。两个同阶拉丁方若叠加后的 $n^2$ 个有序符号对互不相同，称为\keyterm{正交（orthogonal）}。
\end{definition}

对素数 $p$，以模 $p$ 剩余类标记行列，令 $L_a(i,j)=i+aj\pmod p$，其中 $a\ne0$。固定一行时乘以 $a$ 可逆，固定一列时平移可逆，所以每个 $L_a$ 都是拉丁方。若 $a\ne b$，给定叠加符号 $(u,v)$，方程 $i+aj=u,i+bj=v$ 有唯一解 $j=(v-u)/(b-a)$、$i=u-aj$，故 $L_a,L_b$ 正交。这构造了 $p-1$ 个两两正交的拉丁方。

例如 $p=3$ 时，$L_1$ 的三行为 $012,120,201$，$L_2$ 的三行为 $021,102,210$，叠加后九个有序对各出现一次。模合数运算中非零元素未必可逆；例如模 $6$ 乘以 $2$ 会重复余数，不能照搬上述证明。单个拉丁方、两个正交拉丁方以及最大的正交族是不同的存在性问题。

\section{Hamming 距离与编码界}

设字母表 $\mathcal A$ 有 $q\ge2$ 个符号，长度为 $n\ge1$ 的\keyterm{码（code）}是子集 $C\subseteq\mathcal A^n$。字符串 $x,y$ 的\keyterm{Hamming 距离（Hamming distance）} $d_H(x,y)$ 是不同坐标数；对至少含两个码字的 $C$，定义最小距离 $d=\min_{x\ne y\in C}d_H(x,y)$。它满足三角不等式，因为 $x,z$ 不同的坐标必在 $x,y$ 或 $y,z$ 中至少不同一次。

最小距离为 $d$ 的码可以检测至多 $d-1$ 个符号替换错误：这样得到的字符串不可能是另一码字。它能唯一纠正任意至多 $t$ 个替换错误，当且仅当 $2t<d$。充分性来自两个半径 $t$ 的球若相交，会由三角不等式得到码字距离至多 $2t$。必要性来自一对距离为 $d\le2t$ 的码字：在它们不同的坐标中分别选取两方的符号，可构造与两者距离都不超过 $t$ 的接收串，导致歧义。

\begin{theorem}[Hamming 界与 Singleton 界]
令 $t=\lfloor(d-1)/2\rfloor$，则
\begin{equation}
 |C|\sum_{i=0}^t\binom ni(q-1)^i\le q^n,
 \qquad |C|\le q^{n-d+1}.
 \label{eq:coding-bounds}
\end{equation}
\end{theorem}
\begin{proof}
距一个固定码字恰为 $i$ 的字符串通过选择 $i$ 个位置、各换成其他 $q-1$ 个符号之一得到，共 $\binom ni(q-1)^i$ 个。纠错球两两不交，且都包含在 $q^n$ 个字符串中，得到第一式。对第二式，删除任意 $d-1$ 个坐标；若两个码字删除后相同，它们原来的距离至多为 $d-1$，矛盾。因此删坐标映射在 $C$ 上为单射，目标空间有 $q^{n-d+1}$ 个元素。
\end{proof}

二元重复码 $C=\{000,111\}$ 的最小距离为 $3$，可纠正一个替换错误、检测两个错误，两个半径一球恰好覆盖全部八个字符串。若字母表是给定的\keyterm{有限域（finite field）} $\mathbb F_q$，且 $C$ 为 $\mathbb F_q^n$ 的 $k$ 维子空间，则称为 $[n,k,d]_q$ 线性码，含 $q^k$ 个码字。向量的重量是其非零坐标数；线性码中两码字之差仍在码中，所以最小距离等于非零码字的最小重量。例取一个 $3\times7$ 二元矩阵 $H$，其列为七个不同的非零三维向量，令 $C=\{c:Hc^{\mathsf T}=0\}$，称 $H$ 为该码的\keyterm{校验矩阵（parity-check matrix）}。$H$ 包含三个标准基列，秩为三，故 $C$ 维数四；单列非零且两列不同，排除了重量一、二的非零码字，三列 $x,y,x+y$ 又给出重量三码字。因此它是 $[7,4,3]_2$ 码，且 $2^4(1+7)=2^7$，达到半径一的 Hamming 界。

\section{习题}
\begin{enumerate}
  \item\label{ex:designs-codes-1} 设存在 $2$-$(9,3,1)$ 设计，求 $r,b$。仅由这两个整数能否认为已经证明该设计存在？
  \item\label{ex:designs-codes-2} 核对七点设计中每个点所处的区组，说明每两条线也恰好交于一个点。
  \item\label{ex:designs-codes-3} 对 $p=5$ 构造两个正交拉丁方，不必列完表格，但须证明每个有序符号对只出现一次。
  \item\label{ex:designs-codes-4} 若二元长度为 $7$ 的码可纠正任意一次替换错误，利用 Hamming 界求其大小上界，并给出本章中达到上界的例子。
\end{enumerate}

\chapter{组合算法与综合问题}
\label{ch:algorithms}

计数公式、递推和存在性证明可以导出算法，但必须说明算法实际输出什么。输出一个对象、对象总数或全部对象，其资源需求可能相差很大。本章默认输入为有限离散数据，先按算术操作数分析成本，再说明整数位长与输出规模的影响。

\section{枚举、排名与反排名}

枚举长度为 $n$ 的零一串，可以逐位选择 $0$ 或 $1$；枚举排列则逐位选择一个尚未使用的标签。每个合法对象都有唯一的选择路径，所以无重复也无遗漏；每次增加已填位置数，深度至多为 $n$，故过程终止。输出全部 $n!$ 个排列、每个明确写出 $n$ 个元素时，仅写入输出就需要至少与 $n\cdot n!$ 同阶的操作，不能期待关于 $n$ 的多项式总时间。

固定从零开始的字典序编号，\keyterm{排名（ranking）}将对象映射到其编号，\keyterm{反排名（unranking）}从编号恢复对象。对排列 $\pi=(\pi_1,\ldots,\pi_n)$，令 $c_i$ 为尚未使用且小于 $\pi_i$ 的标签数，则
\[
 \operatorname{rank}(\pi)=\sum_{i=1}^n c_i(n-i)!.
\]
因为在第 $i$ 位第一次选择一个较小标签时，后面有 $(n-i)!$ 种补全，且不同首次分歧互斥。逆向给定 $0\le R<n!$，在第 $i$ 步取剩余有序标签中从零编号为 $\lfloor R/(n-i)!\rfloor$ 的元素，再令 $R$ 为相应余数，即可恢复排列。每次商都在合法范围内，商余分解也说明二者互逆。例如 $3142$ 的 $c_i$ 为 $2,0,1,0$，排名为 $13$。

\section{动态规划与状态压缩}

\keyterm{动态规划（dynamic programming）}将重复出现的子问题只计算一次。设不同正整数面值为 $a_1,\ldots,a_k$，目标金额为 $T\ge0$，各面值不限枚数、同面值硬币不可区分。令 $D_j(t)$ 计数只用前 $j$ 种面值凑成 $t$ 的方案，规定 $D_0(0)=1,D_0(t)=0$（$t>0$），并将负金额项置零。按是否使用第 $j$ 种硬币分类，得到
\[
 D_j(t)=D_{j-1}(t)+D_j(t-a_j).
\]
后一项删除一枚第 $j$ 种硬币后仍允许使用该种，因此是 $D_j$ 而非 $D_{j-1}$。一维实现如下：
\begin{verbatim}
CountChange(a[1..k], T):
    dp[0..T] = 0
    dp[0] = 1
    for j = 1..k:
        for t = a[j]..T in increasing order:
            dp[t] = dp[t] + dp[t-a[j]]
    return dp[T]
\end{verbatim}
外层循环不变式是：处理完第 $j$ 种面值，数组为 $D_j$；内层递增保证新项使用当前层的较小金额。算法用 $O(kT)$ 次加法和 $O(T)$ 个整数存储；若 $T$ 以二进制给出，这对输入位长不一定是多项式时间。把金额作为外循环、面值作为内循环则按最后一枚分类，计数的是有序投币串；例如面值 $1,2$ 凑成 $3$，无序方案为两种，有序串为三种。

对 $n$ 人与 $n$ 个不同任务的受限分配，令第 $i$ 人可选任务集为 $A_i\subseteq[n]$，状态 $F(S)$ 表示把前 $|S|$ 人分配到任务集 $S$ 的方案数。初值为 $F(\varnothing)=1$，其余初始为零；按集合大小递增，对每个 $j\in A_{|S|+1}\setminus S$ 更新 $F(S\cup\{j\})\mathrel{+}=F(S)$。最后输出 $F([n])$。删除最后一人的任务唯一确定前驱，所以转移不重不漏；有 $2^n$ 个状态，每态至多 $n$ 条转移，需 $O(n2^n)$ 次加法和 $O(2^n)$ 个整数。答案至多为 $n!$，单个数可能需要 $O(n\log n)$ 位，这一位成本应与操作数分开。

\section{回溯、剪枝与对称消除}

\keyterm{回溯（backtracking）}沿部分解的搜索树逐步扩展；所有合法分支都被访问时保证完备。剪枝必须有“当前前缀没有合法补全”的证明。例如受限分配的剩余二分图若不能饱和剩余人员，则整个分支不可能成功；可以用增广路算法判定，而仅观察每个人都有候选任务是不够的。

若只需对称等价类，可为完整对象建立编码，并保留其轨道中编码字典序最小者。有限轨道总有唯一最小编码，所以每个等价类保留一次；代价是对每个对象检查变换。更高效的部分对称消除必须证明每个轨道仍保留适当数量的代表，不能随意固定某个位置后就把全部计数机械乘回群大小。

直接枚举适合核对小参数，也能发现漏项、边界和重数错误；但有限范围内全部吻合，仍不能证明一般公式。反之，若枚举结果不同，应先核对双方的对象、标签、顺序和对称约定，再判断是哪一个推导出错。

\section{综合专题与方法选择}

对 $x_1+x_2+x_3=8$、$0\le x_i\le4$，容斥给出 $45-3\cdot10=15$；生成函数给出 $[x^8](1+x+\cdots+x^4)^3$；动态规划令 $F_j(t)=\sum_{u=0}^4F_{j-1}(t-u)$，初值 $F_0(0)=1$，同样得 $F_3(8)=15$。容斥适合坏事件交集容易计数时，生成函数适合独立数量限制，动态规划则直接给出可计算的整数表。

六个圆周位置恰有两个红色，其余蓝色，识别旋转后可以按两红点沿圆周的较短距离 $1,2,3$ 分类，恰有三类。Burnside 计算中，恒等旋转固定 $\binom62=15$ 个图案，半周旋转固定三种相对红点，其余旋转不可能保持恰好两红，故平均为 $(15+3)/6=3$。枚举所有十五个二元子集并取旋转规范代表，可以作为独立核对。

对三人的候选集 $\{a,b\},\{a,b\},\{b,c\}$，Hall 定理回答可行性，最大匹配给出一个分配，子集动态规划则给出方案总数为二：第三人必须取 $c$，前两人交换 $a,b$。最大匹配大小本身并不等于匹配方案数。先确定问题输出，才能选对算法。

\section{习题}
\begin{enumerate}
  \item\label{ex:algorithms-1} 计算排列 $2413$ 的零起点字典序排名，并从编号 $13$ 反排名恢复四元素排列。
  \item\label{ex:algorithms-2} 写出不含连续三个 $1$ 的零一串计数算法，给出状态、初值、循环顺序、终止与复杂度。
  \item\label{ex:algorithms-3} 用 $A_1=\{1,2\},A_2=\{2,3\},A_3=\{1,3\}$ 建立子集动态规划，求完整分配数。
  \item\label{ex:algorithms-4} 选择一个小规模计数问题，分别用独立枚举和一种公式方法求解；说明二者相符验证了哪些内容，仍未证明哪些内容。
\end{enumerate}

\appendix

\chapter{预备知识与符号索引}
\label{app:prerequisites}

\section{逻辑、集合与有限求和}

命题 $P\Rightarrow Q$ 只说明 $P$ 足以推出 $Q$，并不自动给出逆命题。等价 $P\Leftrightarrow Q$ 需要两个方向。否定“对每个 $x$ 都存在 $y$ 使 $P(x,y)$ 成立”，应得到“存在 $x$ 使对每个 $y$ 都不成立”；量词顺序决定 $y$ 是否允许依赖 $x$。证明计数公式时，参数范围属于结论的一部分，验证有限个取值不能替代全称命题。

有限不交并的基数可相加；一般并集先变成不交分类或使用容斥。有限和换序只是在重排同一有限项集，例如 $\sum_{k=0}^n\sum_{j=0}^k a_{j,k}=\sum_{j=0}^n\sum_{k=j}^n a_{j,k}$。形式幂级数中的无限表达式则逐系数解释；当每个系数只涉及有限项时，才可以使用本书相应的换序论证。解析无穷级数需要另外的收敛条件。

归纳证明须给出起点、归纳假设和由较小参数向较大参数的推导。强归纳可一次使用所有较小情形；良序原理保证非空非负整数集合有最小元，因此最小反例法与归纳法相容。对递归算法，还须给出严格减小的非负整数参数以证明终止。

\section{行列式与 Cauchy--Binet 公式}

本书线性代数部分使用矩阵乘法、行列式的多线性与交替性，以及秩不超过行列数等基本性质。特别地，线性相关的列使行列式为零，沿仅有一个非零元的行展开可降低阶数。递推求解中使用的 Vandermonde 矩阵为 $V_{ij}=\lambda_j^{i-1}$（$1\le i,j\le r$），其行列式为 $\prod_{i<j}(\lambda_j-\lambda_i)$：行列式在两参数相同时为零，所以被每个差因子整除；双方总次数相同，比较单项式 $\lambda_2\lambda_3^2\cdots\lambda_r^{r-1}$ 的系数确定常数为一。因此不同参数保证可逆。以下公式补足矩阵树定理所需的另一项工具。

\begin{theorem}[Cauchy--Binet 公式]
若 $U$ 为 $r\times m$ 矩阵、$V$ 为 $m\times r$ 矩阵且 $1\le r\le m$，则
\[
 \det(UV)=\sum_{\substack{S\subseteq[m]\\|S|=r}}
             \det(U_{:,S})\det(V_{S,:}),
\]
其中每个子集 $S$ 的指标均按递增顺序选取。
\end{theorem}
\begin{proof}
$UV$ 的第 $j$ 列是 $\sum_{i=1}^m U_{:,i}V_{ij}$。按列多线性展开行列式，得到对 $i_1,\ldots,i_r$ 的求和，项为 $\det(U_{:,i_1},\ldots,U_{:,i_r})\prod_{j=1}^r V_{i_jj}$。若指标重复，行列式为零。对不重复指标按其集合 $S=\{s_1<\cdots<s_r\}$ 归组，写 $i_j=s_{\pi(j)}$，行列式变为 $\operatorname{sgn}(\pi)\det(U_{:,S})$。对排列 $\pi$ 求和的另一因子正是 $\det(V_{S,:})$，于是得到公式。
\end{proof}

\section{概率、符号与模型对照}

对有限样本空间，条件概率在 $\mathbb P(B)>0$ 时定义为 $\mathbb P(A\mid B)=\mathbb P(A\cap B)/\mathbb P(B)$。给定已选择的部分信息后，条件期望就是在剩余可能结果上按更新后的概率加权平均；把下一步结果分类再求平均，可恢复当前条件期望。这是第\ref{ch:probabilistic}章逐步选择方法的依据。

常用计数序列中，$D_n$ 计数错排，$S(n,k)$ 计数 $k$ 块集合分拆，$c(n,k)$ 计数 $k$ 循环置换，$B_n$ 计数全部集合分拆，$C_n$ 计数半长 $n$ 的 Dyck 路径，$p(n)$ 计数整数分拆。其空对象初值为 $D_0=B_0=C_0=p(0)=S(0,0)=c(0,0)=1$；其他边界由具体定义决定。

整数合成保留部件顺序，整数分拆只保留部件重数，集合分拆还保留块内元素标签。标号树保留每个顶点名称，无标号树按图同构取商，平面树另保留孩子次序；它们的计数不应通过未经证明的统一除法互换。普通生成函数的系数直接给对象数，指数生成函数的系数须乘相应阶乘；两者应根据组合运算而非公式外观选用。

\chapter{代表性习题解答}
\label{app:exercises}

本附录选择各章两道习题，给出关键论证或完整计算。未收录的练习应先按其计数对象、边界条件和目标结论独立尝试；计算答案相同并不意味着两种建模必然等价。

\section{基本方法与精确计数}

\subsection{第\ref{ch:objects}章}

习题\ref{ex:objects-1}：先选三名代表再选召集人，有 $\binom53\cdot3=30$ 种；先选召集人，再从其余四人中选两名普通代表，有 $5\binom42=30$ 种。两种过程都唯一指定同一个“委员会加负责人”对象。

习题\ref{ex:objects-2}：纤维两两不交且覆盖 $X$，故 $|X|=\sum_{y\in Y}|f^{-1}(\{y\})|\le d|Y|$。等号当且仅当每个纤维均恰有 $d$ 个元素；若存在一个严格小于 $d$ 的纤维，其余项不超过 $d$，总和便严格小于。

\subsection{第\ref{ch:elementary-counting}章}

习题\ref{ex:elementary-counting-1}：$\mathrm{MISSISSIPPI}$ 中 $I,S$ 各四个，$P$ 两个，$M$ 一个，故不同排列数为 $11!/(4!4!2!)=34650$。每个无副本标签的排列恰对应 $4!4!2!$ 个临时标号排列。

习题\ref{ex:elementary-counting-3}：令 $y_i=x_i-2$，得到四个非负整数之和为 $4$，故方案数为 $\binom73=35$。平移是双射，不改变各盒名称。

\subsection{第\ref{ch:identities}章}

习题\ref{ex:identities-2}：对 $n\ge2$，分解 $k^2=k(k-1)+k$，得到 $n(n-1)2^{n-2}+n2^{n-1}=n(n+1)2^{n-2}$。$n=0$ 时为零，$n=1$ 时为一；后一值也符合上述最终表达式。

习题\ref{ex:identities-3}：分成三段，路径数为 $\binom21\binom32\binom42=36$。若两个必经点互不可比，即一个横坐标更大但纵坐标更小，只向右或向上移动的路径不可能同时经过它们，计数为零。

\subsection{第\ref{ch:pigeonhole}章}

习题\ref{ex:pigeonhole-3}：取序列 $3,2,1,6,5,4,9,8,7$。递增子序列每个递减块至多取一项，所以长度至多三；递减子序列不能跨越数值整体增大的块，也至多长三。

习题\ref{ex:pigeonhole-4}：原操作保持所有数的总和，所以最后是 $1+\cdots+10=55$。改成 $a+b+1$ 后每步总和增加一，而数的个数从十减至一需要九步，最终为 $64$；两种结论都不依赖选取次序。

\subsection{第\ref{ch:inclusion-exclusion}章}

习题\ref{ex:inclusion-exclusion-1}：由错排递推或容斥得 $D_6=265$。恰有两个固定点时，先选固定位置，再让其余四个位置错排，得 $\binom62D_4=15\cdot9=135$。

习题\ref{ex:inclusion-exclusion-3}：容斥得到 $\binom{13}{3}-4\binom93+6\binom53=10$。也可令 $y_i=3-x_i$，得到 $\sum_i y_i=2$；各 $y_i\le3$ 自动满足，隔板法给出 $\binom53=10$。第二种方法利用了总数接近上界之和的特殊结构。

\subsection{第\ref{ch:recurrences}章}

习题\ref{ex:recurrences-3}：令 $a_0=1,a_1=2,a_2=4$。对 $n\ge3$，合法串按尾段分成以 $0$、$01$、$011$ 结尾的三类，删去该尾段分别得到长度 $n-1,n-2,n-3$ 的合法串。因此 $a_n=a_{n-1}+a_{n-2}+a_{n-3}$，继而 $a_3=7,a_4=13,a_5=24$。

习题\ref{ex:recurrences-4}：特征多项式为 $(\lambda-2)^2$，设 $a_n=(A+Bn)2^n$。初值得 $A=1,A+B=2$，故 $a_n=(n+1)2^n$。代入递推两侧均得 $(n+1)2^n$，完成核对。

\section{生成函数与经典结构}

\subsection{第\ref{ch:ogf}章}

习题\ref{ex:ogf-1}：生成函数为 $\frac1{1-x^2}\frac{x}{1-x}(1+x+x^2)$。固定 $x_3=0,1,2$，在 $n=5$ 时分别有 $3,2,2$ 个可选的非负偶数 $x_1$，且 $x_2$ 随即确定为正整数，故共七个解。

习题\ref{ex:ogf-3}：设三种面值的枚数为 $a,b,c$，则 $a+b+c=4,a+2b+5c=10$，相减得 $b+4c=6$。非负条件只允许 $c=1,b=2,a=1$，故只有一种。相应系数为 $[x^{10}y^4]((1-yx)(1-yx^2)(1-yx^5))^{-1}$。

\subsection{第\ref{ch:egf}章}

习题\ref{ex:egf-1}：递推得到 $S(5,2)=15,S(5,3)=25$；两个有名称非空盒子的分配数为 $2^5-2=30$，除去块的两种命名顺序，得到 $S(5,2)=15$。

习题\ref{ex:egf-3}：组件大小只许为正偶数，生成函数为 $\exp(\sum_{j\ge1}x^{2j}/(2j)!)$。四个标签或组成一个四元块，或组成两个二元块，方案数为 $1+\binom42/2=4$。

\subsection{第\ref{ch:classical-structures}章}

习题\ref{ex:classical-structures-2}：凸 $(n+2)$ 边形对应 $C_n$，故七边形对应 $n=5$，三角剖分数为 $C_5=\binom{10}{5}/6=42$。

习题\ref{ex:classical-structures-4}：将 $8$ 拆成八个 $1$，将 $6$ 拆成两个 $3$，再保留原来的 $3$ 和 $1$，得到三个 $3$ 与九个 $1$。总和仍为 $18$，重数的二进制合并又可恢复原分拆。

\section{图、对称与偏序}

\subsection{第\ref{ch:graphs-trees}章}

习题\ref{ex:graphs-trees-2}：依次选出的最小未出现在剩余序列中的标签为 $1,3,2$，连边为 $12,23,24$，最后连 $45$。因此边集为 $\{12,23,24,45\}$，度数依次为 $(1,3,1,2,1)$。

习题\ref{ex:graphs-trees-3}：Pr\"ufer 序列长度为三，其中标签 $1$ 出现两次、$2$ 出现一次，其他不出现，故有 $3!/(2!1!)=3$ 棵树。

\subsection{第\ref{ch:matching-coloring}章}

习题\ref{ex:matching-coloring-3}：四圈的染色多项式为 $(q-1)^4+(q-1)$，所以 $q=3$ 时有 $16+2=18$ 种。颜色已有名称，不再除以颜色排列数。

习题\ref{ex:matching-coloring-4}：顶点数至少三时，边数上界给出平均度数至多 $6-12/n<6$，故有度数至多五的顶点；一、二点图直接成立。删除该顶点后仍可平面，对剩余图按顶点数归纳作六色染色；加回时其至多五个邻点最多禁用五种颜色，尚有一种可选。这证明六色上界，不声称最少颜色数总为六。

\subsection{第\ref{ch:symmetry}章}

习题\ref{ex:symmetry-1}：五点二色旋转固定点总和为 $2^5+4\cdot2=40$，除以五得八类。五个反射各固定 $2^3=8$ 种图案，全部十个变换平均为 $(40+40)/10=8$，所以此例加入反射后数量不变。

习题\ref{ex:symmetry-3}：恒等旋转固定十五个染色，半周旋转固定三对相对红点，其他旋转没有合法固定染色，故得到三类。也可直接按两红点的较短圆周距离为一、二、三分类。

\subsection{第\ref{ch:posets}章}

习题\ref{ex:posets-1}：极小元为 $2,3$，没有最小元；极大元与最大元都为 $6$。最小元必须同时整除集合中的每个元素，不能把任一极小元都称为最小元。

习题\ref{ex:posets-4}：可分为链 $1<2<4<12$ 与 $3<6$（此处小于表示严格整除）。$\{4,6\}$ 为二元反链，任何链划分至少两条；上述构造达到下界，故最少为二。

\section{极值、概率、设计与算法}

\subsection{第\ref{ch:extremal-ramsey}章}

习题\ref{ex:extremal-ramsey-1}：上界为 $\lfloor49/4\rfloor=12$，完全二分图 $K_{3,4}$ 达到它。

习题\ref{ex:extremal-ramsey-3}：$R(3,4)\le R(2,4)+R(3,3)=4+6=10$。这只证明十点足够；要证明最小值恰为十，还需在九点上构造不含红三角形且不含蓝四点团的染色，而递推并未提供这样的构造。

\subsection{第\ref{ch:probabilistic}章}

习题\ref{ex:probabilistic-1}：每个位置成为固定点的概率为 $1/n$，所以 $\mathbb E X=\sum_{i=1}^n1/n=1$。第一矩法的充分条件要求严格小于一，此处不满足；尤其 $n=1$ 时确实没有错排。错排存在性需要别的论证或更细的信息。

习题\ref{ex:probabilistic-3}：函数 $f(p)=np-mp^2$ 是开口向下的二次函数，无约束最大点为 $p=n/(2m)$。若 $2m\ge n$，最大下界为 $n^2/(4m)$；若 $2m<n$，在区间 $[0,1]$ 上最大点为 $p=1$，下界为 $n-m$。实际独立集大小是整数，可将所得实数下界向上取整。

\subsection{第\ref{ch:designs-codes}章}

习题\ref{ex:designs-codes-1}：$r=\lambda(v-1)/(k-1)=4$，$b=vr/k=12$。这只检查了两项必要计数恒等式；若要证明存在，还须给出区组或引用满足全部条件的构造定理。

习题\ref{ex:designs-codes-4}：半径一球大小为 $1+7=8$，所以 $|C|\le2^7/8=16$。由七个不同非零三维二元列组成校验矩阵的零空间给出大小十六、最小距离三的线性码，恰好达到上界。

\subsection{第\ref{ch:algorithms}章}

习题\ref{ex:algorithms-1}：$2413$ 的各位较小未用标签数为 $1,2,0,0$，故排名为 $1\cdot3!+2\cdot2!=10$。编号十三依次给出商 $2,0,1,0$，在剩余有序标签中取对应位置，恢复 $3142$。

习题\ref{ex:algorithms-3}：若第一人取任务一，第二人只能取二，第三人取三；若第一人取二，第二人只能取三，第三人取一。因此完整分配为 $(1,2,3)$ 与 $(2,3,1)$ 两种。子集动态规划在最终状态 $\{1,2,3\}$ 处应得到二，小规模列举在此提供独立核对。

\backmatter
\begin{thebibliography}{9}
\bibitem{mit-combinatorics}
Alexander Postnikov. \emph{18.211 Combinatorial Analysis}, MIT, Fall 2019. 课程主题与教学资料：\url{https://math.mit.edu/~apost/courses/18.211_2019/}.

\bibitem{bogart}
Kenneth P. Bogart. \emph{Combinatorics Through Guided Discovery}. 开放教材：\url{https://bogart.openmathbooks.org/ctgd/ctgd.html}.

\bibitem{dartmouth-combinatorics}
Sergi Elizalde. \emph{Math 68: Algebraic Combinatorics}, Dartmouth College, Fall 2025. 课程主题与讲义入口：\url{https://math.dartmouth.edu/~m68f25/}.

\bibitem{matrix-tree-notes}
MIT OpenCourseWare. \emph{The Matrix-Tree Theorem}. 18.314 Combinatorial Analysis, Fall 2014. \url{https://ocw.mit.edu/courses/18-314-combinatorial-analysis-fall-2014/2724112ea36679f82dc04f0b2f4f355e_MIT18_314F14_mt.pdf}.

\bibitem{zhao-graph-notes}
Yufei Zhao. \emph{Graph Theory and Additive Combinatorics}. MIT 18.225, Fall 2023, Chapter 1. \url{https://ocw.mit.edu/courses/18-225-graph-theory-and-additive-combinatorics-fall-2023/mit18_225_f23_lec_full.pdf}.

\end{thebibliography}

\end{document}
